\section{Design decision log}
\label{sec:decisions}

Append-only log. Newest entries at the bottom. If a decision is reversed, add a
new entry that supersedes the old id; do not silently delete history.

%----------------------------------------------------------------------------
\subsection*{D-001 --- Package rename (structure)}
\begin{decisionbox}[title={D-001 \quad 2026-07-28}]
\textbf{Decision:} Rename/split ambiguous folders to match responsibility:
\texttt{Core$\rightarrow$Game}, \texttt{System$\rightarrow$\{App,Engine,Services\}},
\texttt{Init+Util$\rightarrow$Foundation}, \texttt{AssetLoaders$\rightarrow$Assets}.\\
\textbf{Why:} Names were lying; slowed reasoning and migrations.\\
\textbf{Consequences:} Include/CMake paths updated; archive for old States.
\end{decisionbox}

%----------------------------------------------------------------------------
\subsection*{D-002 --- MathTypes header name}
\begin{decisionbox}[title={D-002 \quad 2026-07-28}]
\textbf{Decision:} Project math aliases live in \pathref{MathTypes.h}, not
\pathref{Math.h}.\\
\textbf{Why:} Windows case-insensitive include collision with CRT
\texttt{math.h}.\\
\textbf{Consequences:} All includes updated; the Foundation package map under
\pathref{docs/packages/} notes the trap.
\end{decisionbox}

%----------------------------------------------------------------------------
\subsection*{D-003 --- Render abstraction direction}
\begin{decisionbox}[title={D-003 \quad 2026-07-28 (implemented; see D-R*) }]
\textbf{Decision:} Command/token submission via \symbolref{IBackend}
+ \symbolref{RenderCommand} queue; game code must not permanently depend on
raw GL for draw submission.\\
\textbf{Why:} Multi-API optionality; clearer frame boundary; testability.\\
\textbf{Consequences:} Phases 1--6 complete for planned scope
(\cref{sec:migration-phases}); hybrid immediate \texttt{Draw} remains only as
fallback.
\end{decisionbox}

%----------------------------------------------------------------------------
\subsection*{D-004 --- Module frame hooks}
\begin{decisionbox}[title={D-004 \quad inferred from IModule}]
\textbf{Decision:} Modules expose \symbolref{Start}/\symbolref{Update}/
\symbolref{DispatchDrawables}/\symbolref{Exit}.\\
\textbf{Why:} Separate simulation tick from draw contribution.\\
\textbf{Consequences:} Draw list is rebuilt or contributed each frame through
dispatch. Module ownership/composition was later closed by D-E1.
\end{decisionbox}

%----------------------------------------------------------------------------
\subsection*{D-005 --- Rules as polymorphic sets}
\begin{decisionbox}[title={D-005 \quad inferred}]
\textbf{Decision:} Each CA variant is a \symbolref{RuleSet} subclass.\\
\textbf{Why:} Open for extension; shared generation sweep.\\
\textbf{Consequences:} Factory logic currently in \symbolref{CellContext}
string matching (could become a registry later).
\end{decisionbox}

%----------------------------------------------------------------------------
\subsection*{D-006 --- Mode handling (enum vs State classes)}
\begin{decisionbox}[title={D-006 \quad current code; author may supersede}]
\textbf{Decision (current):} \symbolref{CellState} enum inside
\symbolref{CellGameModule}.\\
\textbf{Supersedes:} class hierarchy under \pathref{archive/old-states/}.\\
\textbf{Why (inferred):} fewer types for a small mode set.\\
\textbf{Open at decision time:} SAVE/LOAD completeness; whether splash/global
state is acceptable.\\
\textbf{Resolution (2026-08-04):} frame modes are now
\texttt{NORMAL | EDIT | EXIT}; save/load are registered commands with validated
file handling (D-CLI1).
\end{decisionbox}

%----------------------------------------------------------------------------
\subsection*{D-007 --- Enrollment vs per-frame commands}
\begin{decisionbox}[title={D-007 \quad visible in Renderer comments}]
\textbf{Decision:} Resource creation goes through enroll/Create* APIs; do not
stuff mesh/shader creation into the per-frame token stream.\\
\textbf{Why:} Keep the frame queue about \emph{binding and drawing}.\\
\textbf{Consequences:} Need an explicit destroy/reload story for hot-reload
later.
\end{decisionbox}

%----------------------------------------------------------------------------
\subsection*{D-008 --- contributing constraints}
\begin{decisionbox}[title={D-008 \quad from docs/contributing.md}]
\textbf{Decision:} Avoid \texttt{auto}; avoid namespaces (prefer static
classes/structs); no recursion; third-party via PR.\\
\textbf{Why:} Author house style / learning constraints.\\
\textbf{Consequences:} Documentation and new code should respect these unless
explicitly waived in a later decision.
\end{decisionbox}

%----------------------------------------------------------------------------
\subsection*{D-R1 --- RenderCommand payload shape}
\begin{decisionbox}[title={D-R1 \quad 2026-07-28}]
\textbf{Decision:} \symbolref{RenderCommand} uses a \textbf{simple tagged union}
(\texttt{CommandType} + POD payload union), not a kitchen-sink flat struct and
not a byte-stream compiler. Prefer C-style union over \texttt{std::variant}.\\
\textbf{Why:} Type-safe uniforms/clear/update data while remaining
copyable for \symbolref{ArrayQueue}; matches house style.\\
\textbf{Consequences:} \pathref{Rendering/RenderCommand.h} redesigned; execute
path switches on type and reads the matching payload.
\end{decisionbox}

%----------------------------------------------------------------------------
\subsection*{D-R2 --- Who emits tokens}
\begin{decisionbox}[title={D-R2 \quad 2026-07-28}]
\textbf{Decision:} Each drawable emits its own tokens via
\texttt{AppendCommands(Renderer*)} (or equivalent). \symbolref{Renderer}
owns frame setup (clear/viewport) and submit.\\
\textbf{Why:} Smallest strangler step from today's
\symbolref{DispatchDrawables} list.\\
\textbf{Consequences:} Hybrid phase may still call immediate \texttt{Draw}
for unmigrated drawables after token submit.
\end{decisionbox}

%----------------------------------------------------------------------------
\subsection*{D-R3 --- Handle typing for v1}
\begin{decisionbox}[title={D-R3 \quad 2026-07-28}]
\textbf{Decision:} Keep \texttt{unsigned long} opaque handles with backend
registries (\texttt{tableID} $\rightarrow$ GL object). Strong typedefs deferred.\\
\textbf{Why:} Minimal churn; critical fix is resolve-on-submit, not type names.\\
\textbf{Consequences:} Execute must never treat handles as raw GLuints without
registry lookup.
\textbf{Status:} Historical v1 decision; superseded by D-R16.
\end{decisionbox}

%----------------------------------------------------------------------------
\subsection*{D-R4 --- Migration scope}
\begin{decisionbox}[title={D-R4 \quad 2026-07-28}]
\textbf{Decision:} Full migration Phases 1--5 (token proof $\rightarrow$ Scene
routing $\rightarrow$ Canvas $\rightarrow$ UI drawables $\rightarrow$ cleanup).
Phase 6 (mock backend) included; second real GPU API still optional.\\
\textbf{Why:} Playable strangler path to API isolation without big-bang rewrite.\\
\textbf{Also freeze for v1:} destroy resources at shutdown only; Windows GL is
primary platform. The lifetime clause is superseded by D-R16/D-R17.\\
\textbf{Superseded in part by D-P1/D-P4:} full-frame Canvas rewrite was v1 only.
Dirty-rect/PBO uploads remain, while the R8 presentation was later superseded
by the restored RGB fade path documented in D-P4.
\end{decisionbox}

%----------------------------------------------------------------------------
\subsection*{D-R5 --- Blend state on first enable}
\begin{decisionbox}[title={D-R5 \quad 2026-07-28}]
\textbf{Decision:} When enabling blending in \symbolref{GLDevice::ApplyPipelineState},
always call \texttt{glBlendFunc} (do not rely on CPU-side factor cache matching
defaults).\\
\textbf{Why:} OpenGL default is \texttt{ONE}/\texttt{ZERO}; cache defaults were
SrcAlpha/OneMinusSrcAlpha so factors looked ``unchanged'' and console panel
lost transparency after the token migration.\\
\textbf{Consequences:} Semi-transparent UI (console panel alpha) works again.
\end{decisionbox}

%----------------------------------------------------------------------------
\subsection*{D-R6 --- Dead stack cleanup (Phase 5)}
\begin{decisionbox}[title={D-R6 \quad 2026-07-28}]
\textbf{Decision:} Archive window \symbolref{RenderQueue}, legacy
\texttt{Drawable.cpp} GL helpers, empty \texttt{Transform}/\texttt{RenderContext},
and unused \texttt{CellCommandLine} under \pathref{archive/dead-render/}.
\symbolref{DrawableBase} is a pure list/token interface (no GL handles).\\
\textbf{Why:} Token path won; dual submission models confuse maintenance.\\
\textbf{Consequences:} \texttt{IRenderWindow::getRenderQueue} removed; frame
draw list is only \symbolref{Scene} + \symbolref{Renderer}.
\end{decisionbox}

%----------------------------------------------------------------------------
\subsection*{D-R7 --- Mock backend for tests (Phase 6)}
\begin{decisionbox}[title={D-R7 \quad 2026-07-28}]
\textbf{Decision:} Add headless \texttt{MockBackend} implementing
\symbolref{IBackend}: record \texttt{Create*} and snapshot the command queue
on \texttt{SubmitCommandQueue}. Originally shipped as CMake target
\texttt{IllumoRenderTests} / CTest \texttt{MockBackend}; the current library
runner is \texttt{IllumoTests}, with exact \texttt{Illumo.*} CTest cases. The
old target and single-test names are no longer retained. Do \textbf{not} add
Vulkan/Metal yet.\\
\textbf{Why:} First meaningful automated render test without a GPU context;
proves token vocabulary and handle plumbing independently of GL.\\
\textbf{Consequences:} Library tests use \texttt{Illumo::TestSupport}; workspace
CI aggregates the library and game runners under the \texttt{IllumoWorkspace}
label. A second graphics backend remains future work after production needs
appear.
\end{decisionbox}

%----------------------------------------------------------------------------
\subsection*{D-R8 --- Inject IBackend into Renderer}
\begin{decisionbox}[title={D-R8 \quad 2026-07-28, refined 2026-08-07}]
\textbf{Decision:} \symbolref{Renderer} is backend-neutral: its only constructor
takes an injected \symbolref{IBackend*} and a \texttt{takeOwnership} flag.
Production library initialization constructs the real backend via
\texttt{CreateOpenGLBackend} under \pathref{Rendering/OpenGL/}, initializes it
exactly once, and transfers a \texttt{std::unique\_ptr<IBackend>} into the
renderer. Tests inject \texttt{MockBackend} through the same ownership
boundary. \texttt{Renderer.h} must not include or construct concrete OpenGL
types.\\
\textbf{Why:} End-to-end drawable $\rightarrow$ enroll/token $\rightarrow$ submit
without a GPU; second platforms can supply a different \symbolref{IBackend}
without touching the generic frame-submit surface.\\
\textbf{Consequences:} \texttt{TestRendererE2E} covers inject, owned-backend
composition style, header neutrality, \texttt{RenderScene}, Canvas tokens,
hybrid immediate drawables, and \texttt{RenderProofQuad}. OpenGL sources stay
under \pathref{Rendering/OpenGL/}; Mock under \pathref{Rendering/Mock/}.
\end{decisionbox}

%----------------------------------------------------------------------------
\subsection*{D-P1 --- Dirty visual path}
\begin{decisionbox}[title={D-P1 \quad 2026-07-28}]
\textbf{Decision:} Canvas/CellGameModule use dirty flags so idle frames skip
full-grid visual work and GPU upload.\\
\textbf{Why:} Dominant cost was recoloring + full texture upload every render
frame even when the CA did not step.\\
\textbf{Consequences:} \texttt{cellsDirty}, upload pending, early-out
\texttt{tickVisual}; still life after sim step does not re-dirty.
\end{decisionbox}

%----------------------------------------------------------------------------
\subsection*{D-P2 --- UI batching and geometry cache}
\begin{decisionbox}[title={D-P2 \quad 2026-07-28}]
\textbf{Decision:} CommandLine emits one packed \texttt{UpdateBuffer}/draw when
open and dirty; GLString caches GPU geometry until content/style changes.
Settled CommandLine frames keep one \texttt{DrawIndexed} and skip the buffer
upload.\\
\textbf{Why:} Per-line uploads and FPS string rebuilds were measurable when the
console/FPS overlay was active. Re-wrapping the full history and baking a
per-frame accent pulse kept the open console expensive.\\
\textbf{Consequences:} Token stream stays small; FPS may rebuild $\sim$1\,Hz;
idle console history reuses the last mesh until scroll, input, layout, caret
phase, or a quantized pulse changes.
\end{decisionbox}

%----------------------------------------------------------------------------
\subsection*{D-P3 --- Double-buffer CA generation}
\begin{decisionbox}[title={D-P3 \quad 2026-07-28}]
\textbf{Decision:} \texttt{RuleSet::calcGeneration} writes pure
\texttt{nextState(cell, aliveNeighbors)} into a scratch buffer, then
\texttt{memcpy}s back; neighbor counts use raw grid pointers + toroidal wrap.\\
\textbf{Why:} Correct simultaneous update; faster than per-cell
\texttt{getCanvasPixel}/in-place virtual writes; sparse dirty AABB for visuals.\\
\textbf{Consequences:} All active rulesets override \texttt{nextState}; stubs
use identity.\\
\textbf{Refined by D-P5 (2026-08-06):} scratch lives on \symbolref{CellGrid}
as a back buffer; promotion is pointer swap, not full-grid \texttt{memcpy}.
\end{decisionbox}

%----------------------------------------------------------------------------
\subsection*{D-P5 --- Single-pass dirty AABB + buffer swap}
\begin{decisionbox}[title={D-P5 \quad 2026-08-06}]
\textbf{Decision:} Track changed-cell AABB during the eval pass; store next
generation in \texttt{CellGrid} back buffer; on any change call
\texttt{swapLifeBuffers()} and mark the sparse dirty region. Interior Moore
counts skip toroidal wrap branches.\\
\textbf{Why:} Second full-grid compare pass and full \texttt{memcpy} write-back
dominated serial generation cost on large dense grids.\\
\textbf{Consequences:} \texttt{lifeCanvas} always means current front; rules
must not retain raw pointers across a swap; still life leaves domain clean.
\end{decisionbox}

%----------------------------------------------------------------------------
\subsection*{D-P6 --- Fade loop + packed dirty PBO staging}
\begin{decisionbox}[title={D-P6 \quad 2026-08-06}]
\textbf{Decision:} Keep CPU RGB fade presentation. Hoist coordinates in
\texttt{tickVisual}/\texttt{snapVisualToTargets}; rectangular target rebuild
uses one \texttt{includeRect}. \texttt{GLTexture::UpdateSubImage} packs partial
dirty rects tightly and maps only the staging range (no full-buffer orphan
every frame).\\
\textbf{Why:} Per-cell div/mod and full-texture PBO orphan were measurable on
large dirty rects while fade quality must stay.\\
\textbf{Consequences:} Fade math unchanged; MockBackend token path unaffected.
\end{decisionbox}

%----------------------------------------------------------------------------
\subsection*{D-P7 --- Optional row-parallel generation}
\begin{decisionbox}[title={D-P7 \quad 2026-08-06}]
\textbf{Decision:} \texttt{RuleSet::calcGeneration} may partition rows across
workers when cell count $\ge 512\times 512$ (auto) or when
\texttt{setWorkerOverride(N)} forces $N$ workers. Results are bit-identical to
serial. Presentation/fade stays main-thread.\\
\textbf{Why:} Large dense grids are eval-bound after D-P5; spawn cost makes
parallel unhelpful on default $80\times 60$.\\
\textbf{Consequences:} Tests cover serial/parallel identity
(\texttt{IllumoGame.Sim.SerialParallelIdentical}); micro-benches under
\texttt{IllumoGame.Sim.MicroBench}. Persistent thread pool deferred.
\end{decisionbox}

%----------------------------------------------------------------------------
\subsection*{D-P4 --- dirty-rect PBO uploads (R8 path partially superseded)}
\begin{decisionbox}[title={D-P4 \quad 2026-07-28}]
\textbf{Decision (original):} Present cells as \texttt{GL\_R8} texture of state
bytes plus a $256\times 1$ RGB palette (\texttt{evalCell}). Upload dirty rects
via PBO ping-pong in \texttt{GLTexture::UpdateSubImage}. Drop dual float RGB
staging.\\
\textbf{Why:} Large grids: bandwidth and CPU of RGB float fade dominated.\\
\textbf{Consequences (kept):} dirty-rect \texttt{UpdateTexture}, PBO path,
bind-state tracker in \symbolref{GLDevice}.\\
\textbf{Superseded presentation (2026-07-30):} live path restored to CPU RGB
fade (\texttt{targetRgb}/\texttt{displayRgb}) + RGB display texture so cell
colors ease rather than snap. R8+palette is historical experiment, not current
code truth. Canonical: \pathref{docs/architecture-consensus.md}.
\end{decisionbox}

%----------------------------------------------------------------------------
\subsection*{D-E1 --- Module registration lives in App}
\begin{decisionbox}[title={D-E1 \quad 2026-07-28}]
\textbf{Decision:} \symbolref{Illumo} does not include or construct
\symbolref{CellGameModule}. Today,
\pathref{IllumoGame/Source/App/CellMain.cpp} applies product configuration,
calls fallible \texttt{initialize}, then registers the required game module and
the optional Debug module before \texttt{startModules}.\\
\textbf{Why:} Engine must not depend on Game headers; product composition
belongs at the application boundary.\\
\textbf{Consequences:} New products/tests can register different module sets;
Engine only knows \symbolref{IModule}.
\end{decisionbox}

%----------------------------------------------------------------------------
\subsection*{D-E2 --- InputManager has no Game types}
\begin{decisionbox}[title={D-E2 \quad 2026-07-28}]
\textbf{Decision:} Remove \texttt{CellContext*} and Game includes from
\symbolref{InputManager}. Also drop unused \texttt{SplashText} extern from the
input header.\\
\textbf{Why:} Services must not depend on Game; input is a pure device/action
service. Splash and rules stay in Game/Rendering.\\
\textbf{Consequences:} InputManager header only needs KeyCode/InputContext/GLFW.
\end{decisionbox}

%----------------------------------------------------------------------------
\subsection*{D-E3 --- EntityTable is not the cell path}
\begin{decisionbox}[title={D-E3 \quad 2026-07-28}]
\textbf{Decision:} Archive \pathref{EntityTable.h} and \pathref{ModuleObject.h}
under \pathref{archive/dead-engine/}. Remove \texttt{entityTable} from
\symbolref{Scene}/\symbolref{IllumoContext}.\\
\textbf{Why:} Nothing allocated or used an EntityTable; the live path is
drawables + tokens, not ECS.\\
\textbf{Consequences:} Reintroduce a general object table only when a real
feature needs it; document that cells are not entities.
\end{decisionbox}

%----------------------------------------------------------------------------
\subsection*{D-E4 --- Scene is a drawable list, not a graph}
\begin{decisionbox}[title={D-E4 \quad 2026-07-28}]
\textbf{Decision:} \symbolref{Scene} holds only \texttt{drawables},
\texttt{window}, and \texttt{activeCamera}. Remove \texttt{root},
\texttt{nodeLookup}, no-op \texttt{Update}, and unused
\texttt{RenderableObject}. Archive \pathref{SceneObject.h}; \symbolref{Camera}
and \symbolref{CommandLine} no longer inherit a node graph.\\
\textbf{Why:} Architecture review (dead graph scaffolding); production never
used parent/child transforms for cells or UI.\\
\textbf{Consequences:} Frame contribution model is explicit and small.
\end{decisionbox}

%----------------------------------------------------------------------------
\subsection*{D-E5 --- Freeze IllumoContext; validate at Start}
\begin{decisionbox}[title={D-E5 \quad 2026-07-28}]
\textbf{Decision:} Keep \symbolref{IllumoContext} as a service bag for the
\textbf{two shipped modules} only. Document required fields per module.
\texttt{Start()} logs an error and bails if the bag is incomplete.
Do not grow the bag for convenience; a third module should take explicit
constructor dependencies if its subset differs.\\
\textbf{Why:} Architecture review --- bag cost is zero at $N=2$ but becomes a
de-facto global API as $N$ grows.\\
\textbf{Consequences:} \texttt{IllumoContextHasGameCore} /
\texttt{IllumoContextHasDebugCore} helpers; freeze note on the struct.
\end{decisionbox}

%----------------------------------------------------------------------------
\subsection*{D-C1 --- Canvas dual role is intentional (for now)}
\begin{decisionbox}[title={D-C1 \quad 2026-07-28}]
\textbf{Decision:} Keep \symbolref{Canvas} as domain grid + fade view + GPU
enroll in one type. Document the scale wall: full-grid
\texttt{calcGeneration}, dual float RGB buffers ($\sim$24\,B/cell), no
chunked simulation domain.\\
\textbf{Why:} Locality for dirty-rect + palette + fade; headless tests already
use injected MockBackend rather than pure LifeGrid.\\
\textbf{Consequences:} Split \texttt{LifeGrid}/\texttt{CanvasView} only if pure
sim tests or large canvases become product requirements.
\end{decisionbox}

%----------------------------------------------------------------------------
\subsection*{D-R9 --- String-named uniforms are backend debt}
\begin{decisionbox}[title={D-R9 \quad 2026-07-28}]
\textbf{Decision:} Keep \texttt{char name[32]} uniforms for OpenGL + Mock.
Treat as debt if a second \emph{real} GPU backend is adopted; prefer binding
points / pre-resolved locations then.\\
\textbf{Why:} Architecture review --- uniform path is the least portable part
of \symbolref{RenderCommand}; handles and UpdateTexture are closer to neutral.\\
\textbf{Consequences:} Comment on uniform payloads; no API change now.
\end{decisionbox}

%----------------------------------------------------------------------------
\subsection*{D-R10 --- Production drawables are pure-token}
\begin{decisionbox}[title={D-R10 \quad 2026-07-28}]
\textbf{Decision:} \symbolref{Canvas}, \symbolref{CommandLine},
\symbolref{GLString}, \symbolref{SplashText} use
\texttt{AppendCommands} only (return true when contributing). Immediate
\texttt{Draw()} hybrid remains for tests/stubs. CRTP \texttt{DrawImpl} is
legacy for that fallback.\\
\textbf{Why:} Migration complete for shipped UI/canvas; dual path is debt only
where unmigrated types exist.\\
\textbf{Consequences:} Documented in \pathref{Drawable.h} and
\texttt{Renderer::RenderScene}.
\end{decisionbox}

%----------------------------------------------------------------------------
\subsection*{D-R14 --- Scene layers and built-in render styles}
\begin{decisionbox}[title={D-R14 \quad 2026-08-07}]
\textbf{Decision:} Introduce \textbf{Phase A} composition without multi-pass
GPU targets:
\begin{itemize}
  \item \symbolref{Scene} buckets drawables into ordered layers
        (\texttt{World}, \texttt{UI}, \texttt{Debug}) within one main pass
        (default framebuffer, single clear).
  \item \symbolref{Renderer} owns built-in \symbolref{RenderStyle} entries
        (\texttt{Canvas}, \texttt{UiText}, \texttt{Console}): each holds a
        shader table handle plus \texttt{PipelineState} defaults.
  \item Production drawables keep content handles (mesh/texture/domain data)
        and call \texttt{bindStyle} before content tokens.
\end{itemize}
Layers are not GPU target/load-store passes. Unimplemented pass scaffolding is
removed; explicit offscreen targets remain deferred.\\
\textbf{Why:} World/UI/debug stacking and shared draw policy should not live
as duplicated per-drawable shader enrollment and ad-hoc pipeline structs.\\
\textbf{Consequences:} \pathref{RenderLayerId.h}, \pathref{RenderStyle.h},
\pathref{RendererStyles.cpp}; modules pass a layer to
\texttt{Scene::AddDrawable}; docs distinguish layer vs pass.
\end{decisionbox}

%----------------------------------------------------------------------------
\subsection*{D-R15 --- Render primitives and GameVisual host}
\begin{decisionbox}[title={D-R15 \quad 2026-08-07}]
\textbf{Decision:} Introduce value-type \textbf{render primitives}
(\symbolref{ShapePrimitive}, \symbolref{SpritePrimitive}) composed on a
\symbolref{GameVisual} \texttt{Drawable} host.
\begin{itemize}
  \item Primitives are not Scene entries; the host is.
  \item Game/editor objects \textbf{embed} \symbolref{GameVisual} (or multiple)
        and add primitives to build complex visuals without hand-rolling
        enroll/bind/draw tokens.
  \item Built-in styles \texttt{Shape} and \texttt{Sprite} live on
        \symbolref{Renderer}; GPU mesh/texture/shader objects stay in the
        backend registries.
  \item Coordinate space is host policy (\texttt{Pixels} or \texttt{World}).
\end{itemize}
\textbf{Why:} Composition over per-primitive drawables keeps Scene traffic low
and matches D-R14 style ownership.\\
\textbf{Consequences:} \pathref{Rendering/Primitives/}; headless tests
\texttt{Illumo.GameVisual.*}; Canvas/UI paths unchanged.
\end{decisionbox}

%----------------------------------------------------------------------------
\subsection*{D-R16 --- Typed resource lifetime and bounded command growth}
\begin{decisionbox}[title={D-R16 \quad 2026-08-09}]
\textbf{Decision:} Replace raw renderer IDs with non-convertible
\texttt{MeshHandle}, \texttt{ShaderHandle}, \texttt{TextureHandle}, and
\texttt{RenderStyleHandle} values containing slot and generation. Backends
allocate handles and validate replace, destroy, validity, metadata, and command
resolution. Stale operations warn and safely no-op. Replace the fixed token
array with vector storage that reserves 2,048, grows as needed, and rejects only
at a configurable 65,536 default ceiling while reporting high-water and
rejection counts. Scissor state is explicitly enabled or disabled.\\
\textbf{Why:} Reusable rendering needs safe lifetime, hot replacement, and
realistic frame growth without silent corruption or unbounded memory.\\
\textbf{Consequences:} D-R3 and the fixed-capacity part of D-R12 are
superseded. Unimplemented command/pass scaffolding is removed. OpenGL and Mock
implement the same lifecycle contract.
\end{decisionbox}

%----------------------------------------------------------------------------
\subsection*{D-R17 --- Managed texture and shader assets}
\begin{decisionbox}[title={D-R17 \quad 2026-08-09}]
\textbf{Decision:} \symbolref{AssetManager} caches file textures by canonical
path plus options and shaders by canonical path pair. Duplicate acquisition
increments a reference count; final release destroys the backend resource.
Async is default, synchronous loading is available for startup/tests, one owned
worker performs file read/decode, and render-thread \texttt{pump} performs GPU
replacement. Stable fallbacks cover pending/failed initial loads. Failed reload
keeps the last good object. Debug polls timestamps every 500\,ms; all builds
retain explicit reload.\\
\textbf{Why:} Future 2D projects need deterministic resource ownership and live
iteration without compiling or mutating GPU objects from a worker.\\
\textbf{Consequences:} Managed scope is textures and shaders only. Dynamic quad
meshes stay renderer-owned; font atlases, model import, and general 3D meshes
remain outside this milestone.
\end{decisionbox}

%----------------------------------------------------------------------------
\subsection*{D-R18 --- Painter-correct 2D primitive stream}
\begin{decisionbox}[title={D-R18 \quad 2026-08-09}]
\textbf{Decision:} Give \symbolref{GameVisual} one stable ordered stream across
shape, sprite, and text values. Integer draw order is followed by insertion
sequence; only adjacent compatible items batch. Parent/local
\texttt{Transform2D}, normalized pivots, atlas regions, flips, and centered
sprite helpers are supported without a hierarchy. Dynamic quad buffers begin at
1,024, double, and stop at a configurable 65,536 default ceiling.
\symbolref{SpriteAnimator} is caller-updated and exposes only current visual
state with Once, Loop, and PingPong modes.\\
\textbf{Why:} Global texture sorting breaks overlapping alpha composition, and
animation time inside submission makes rendering nondeterministic.\\
\textbf{Consequences:} Top-left behavior remains the default $(0,0)$ pivot;
rendering observes state and never owns gameplay time.
\end{decisionbox}

%----------------------------------------------------------------------------
\subsection*{D-R19 --- Product proof before renderer extraction}
\begin{decisionbox}[title={D-R19 \quad 2026-08-09}]
\textbf{Decision:} Illumo is a cellular-automata application with a reusable 2D
renderer, not a general engine. A Debug-only first-party atlas and managed
sprite-shader showcase plus
\texttt{renderer\_demo}/\texttt{assets}/\texttt{asset\_reload} commands prove
the foundation in the product. Extract a standalone library only when a second
real project proves the required public boundary.\\
\textbf{Why:} A real consumer reveals the useful boundary better than detached
ECS, scene-graph, render-graph, or multi-backend scaffolding.\\
\textbf{Consequences:} Release composition remains unchanged. Follow-on work is
font/UI layout, tilemaps/culling/particles, then offscreen composition as actual
projects require.\\
\textbf{Status:} Superseded by D-E6. The approved consumer is the repository's
\texttt{IllumoGame} product, so the extraction records an explicit product
boundary; it is not evidence of a second independent product and does not
authorize speculative general-engine work.
\end{decisionbox}

%----------------------------------------------------------------------------
\subsection*{D-R20 --- Primitive-composed product UI}
\begin{decisionbox}[title={D-R20 \quad 2026-08-09}]
\textbf{Decision:} Build Illumo's existing UI drawables from
\symbolref{GameVisual} shape and text primitives. Keep
\symbolref{CommandLine} as the owner of console layout and interaction, retain
\symbolref{GLString}'s geometry cache while allowing optional panel chrome, and
use that decorated-label path for Debug FPS and fading EDIT/NORMAL notices.
Share colors and compact panel metrics through value-only \texttt{UiTheme}; do
not introduce a retained widget tree, layout engine, or additional render pass.\\
\textbf{Why:} The renderer foundation should prove useful composition in the
real product, and the console, FPS, and mode notice need one coherent visual
language without increasing Scene traffic or creating a second UI architecture.\\
\textbf{Consequences:} Console chrome now uses fills, outline rectangles, and
line primitives instead of a bespoke packed-quad aesthetic. Decorated labels
compose shadow, surface, border, accent, and text in one adjacent batch. Existing
editing, floating-window, fading, Scene-layer, and Debug/Release ownership
boundaries remain unchanged; headless tests assert the primitive composition.
\end{decisionbox}

%----------------------------------------------------------------------------
\subsection*{D-F1 --- MacroDefs Windows include deferred}
\begin{decisionbox}[title={D-F1 \quad 2026-07-30}]
\textbf{Decision:} Do \textbf{not} split \pathref{MacroDefs.h} /
\texttt{Windows.h} now. Revisit only if include toxicity causes real pain
(compile time, macro collisions, portability friction).\\
\textbf{Why:} Author priority --- convenience of \texttt{\_\_SLEEP} outweighs
hygiene until symptoms appear.\\
\textbf{Consequences:} Status note in services/platform section; not an open
blocker.
\end{decisionbox}

%----------------------------------------------------------------------------
\subsection*{D-N1 --- Illumo is the project and product name}
\begin{decisionbox}[title={D-N1 \quad 2026-08-04}]
\textbf{Decision:} Adopt \textbf{Illumo} as the repository-facing project,
product, CMake target, runtime title, console name, and save-format name. Keep
the existing \symbolref{Illumo} runtime-host class; the shared name is
intentional. Canonical tests are \texttt{IllumoTests}, and new saves use the
\texttt{.illumo} extension.\\
\textbf{Why:} The author selected Illumo as the permanent product identity.\\
\textbf{Consequences:} Live first-party source and documentation use Illumo;
generated build trees, vendored material, and archived historical artifacts are
not implementation sources and are not rewritten by this decision.\\
\textbf{Status:} Superseded by D-N2 and D-N4 for executable, product-visible,
test, and new-save filename identity. The Illumo repository/library/host name
and legacy \texttt{.illumo} loading remain valid.
\end{decisionbox}

%----------------------------------------------------------------------------
\subsection*{D-B1 --- DebugModule is excluded from Release}
\begin{decisionbox}[title={D-B1 \quad 2026-08-04}]
\textbf{Decision:} \symbolref{DebugModule} is Debug-build product composition
only. \symbolref{CellMain} must include/register it behind
\texttt{\#ifndef NDEBUG}, and CMake must compile its implementation only for the
Debug configuration.\\
\textbf{Why:} The developer overlay, console input handling, FPS/debug labels,
and instrumentation are development tools, not Release product behavior.\\
\textbf{Consequences:} Shared command/logging service code may remain available
to tests or non-UI systems, but Release must contain no DebugModule object or
module registration. Audit this boundary after CMake/composition changes.
\end{decisionbox}

%----------------------------------------------------------------------------
\subsection*{D-CLI1 --- Split generic and domain console ownership}
\begin{decisionbox}[title={D-CLI1 \quad 2026-08-04}]
\textbf{Decision:} \symbolref{CommandLine} owns editing, parsing, help,
environment/application commands, and rendering. Product modules register
domain behavior through \symbolref{CommandRegistry}; each registration carries
usage, description, and optional completion candidates.\\
\textbf{Why:} Services should not acquire Game dependencies, while help and
completion must describe the commands that actually exist.\\
\textbf{Consequences:} \symbolref{CellGameModule} owns simulation, canvas,
ruleset, camera, and save/load commands. The shared known-ruleset list drives
validation/completion. Queued commands return successfully without falling
through to environment lookup or an unknown-command error.
\end{decisionbox}

%----------------------------------------------------------------------------
\subsection*{D-UI1 --- Measured console editor in one bounded batch}
\begin{decisionbox}[title={D-UI1 \quad 2026-08-04}]
\textbf{Decision:} Render the command caret and selection from measured text
prefixes, keep long input visible through a horizontal viewport, and retain one
batched token draw for console chrome and text. Size the dynamic mesh for 6,000
UI quads and regression-test a detailed help page beyond the former cap.\\
\textbf{Why:} An appended underscore produced an apparent cursor offset, while
the 2,000-quad mesh silently truncated easy-font output because each character
uses several quads.\\
\textbf{Consequences:} Editing supports selection and word operations without
lying about insertion position; a full help page remains one draw without
cutting off its final line.
\end{decisionbox}

%----------------------------------------------------------------------------
\subsection*{D-UI2 --- Wrapped console history with unified scroll metrics}
\begin{decisionbox}[title={D-UI2 \quad 2026-08-06}]
\textbf{Decision:} Word-wrap console history to the panel width, scroll by
visual lines rather than raw history entries, and compute PageUp/wheel/scrollbar
limits from the same floating-aware panel layout used when drawing. Raise the
console UI mesh to 8,000 quads to leave headroom for wrapped lines plus chrome.\\
\textbf{Why:} Hard truncation hid the rest of long help lines, and scroll
handlers used a mounted-only height estimate that disagreed with the floating
panel, so paging to the oldest history could leave text clipped or unreachable.\\
\textbf{Consequences:} Scrolling fully to the start of history shows the oldest
characters completely; floating resize and mounted mode share one max-scroll
definition; wrap metrics are cached until history or panel width change and
must stay equivalent to direct wrapping; a focused wrap/scroll regression
covers the clamp and batch path.
\end{decisionbox}

%----------------------------------------------------------------------------
\subsection*{D-DOC1 --- One first-party documentation tree}
\begin{decisionbox}[title={D-DOC1 \quad 2026-08-04}]
\textbf{Decision:} Keep all first-party project documentation under
\pathref{docs/}: current Markdown at the top level, package maps in
\pathref{packages/}, dated records in \pathref{sessions/}, provenance in
\pathref{history/}, one current LaTeX entrypoint at \pathref{latex/illumo.tex},
and generated output in \pathref{output/}.\\
\textbf{Why:} Competing LaTeX entrypoints and package READMEs scattered beside
code made current truth and PDF generation difficult to discover.\\
\textbf{Consequences:} The repository root keeps only the landing README and
AGENTS bootstrap as intentional prose exceptions. Licenses and vendored notes
remain beside governed assets. Architecture changes update Markdown, relevant
LaTeX chapters, and this append-only log together. \pathref{docs/build.ps1}
is the PDF wrapper and the optional default-build \texttt{IllumoDocs} target
uses that same script when TeX is available.
\end{decisionbox}

%----------------------------------------------------------------------------
\subsection*{D-T1 --- Granular CTest registration and coverage gate}
\begin{decisionbox}[title={D-T1 \quad 2026-08-04}]
\textbf{Decision:} Keep two compile-efficient runners: \texttt{IllumoTests}
for the reusable library and \texttt{IllumoGameTests} for product behavior.
Register each logical behavior as an exact owner-prefixed
\texttt{Illumo.<area>.<case>} or \texttt{IllumoGame.<area>.<case>} CTest entry.
CTest runs one filtered case per process in an isolated working directory and
aggregates both inventories under \texttt{IllumoWorkspace}. With
\texttt{ILLUMO\_ENABLE\_COVERAGE=ON}, Clang/LLVM adds the
\texttt{IllumoCoverage} target and enforces at least 85\% line coverage over
headless-testable first-party production code.\\
\textbf{Why:} A monolithic CTest result hid which of dozens of behaviors ran or
failed, and a narrow denominator could report a strong percentage while omitting
game commands, persistence, input, or assets.\\
\textbf{Consequences:} New behaviors receive separate stable names; file-backed
tests may run in parallel without sharing working files; combined coverage
merges both runners. Tests, vendored/system code, TestSupport, and the live
OpenGL backend are excluded from the automated production denominator. Native
dialogs, real-window behavior, and live OpenGL still require proportional smoke
tests. This supersedes the single-CTest-name registration detail in D-R7
without changing the MockBackend decision.
\end{decisionbox}

%----------------------------------------------------------------------------
\subsection*{D-UI3 --- Floating mode, title-bar dragging, and corner resize}
\begin{decisionbox}[title={D-UI3 \quad 2026-08-06}]
\textbf{Decision:} Extend \symbolref{CommandLine} to support dynamic window modes and resizable floating window geometry. Mode switching is driven by double-clicking the console title bar or executing \texttt{console\_mode [floating|mounted|toggle]}. In floating mode, left-clicking and dragging the title bar repositions the window across the screen, and dragging the bottom-right corner grip handle (or running \texttt{console\_size <W> <H> | reset}) resizes the console window dynamically with real-time viewport clipping and single-batch quad rendering.\\
\textbf{Why:} Enhances Developer Console usability, allows moving and scaling the console UI away from active grid regions, and provides modular window layout flexibility.\\
\textbf{Consequences:} \symbolref{CommandLine} geometry (chrome, history text, input trough, scrollbar, ghost text, selection highlight) is computed relative to dynamic floating bounds \texttt{panelX0}, \texttt{panelX1}, \texttt{panelY0}, \texttt{panelY1}; mouse press/drag event dispatch accounts for floating position offset and corner resize bounds.
\end{decisionbox}

%----------------------------------------------------------------------------
\subsection*{D-UI4 --- Plain developer console with levelled, filterable output}
\begin{decisionbox}[title={D-UI4 \quad 2026-09-23}]
\textbf{Decision:} The developer console stops using the product
\texttt{UiTheme} and draws with its own flat terminal palette: an opaque dark
panel, one-pixel border and rules, a header, an output band anchored at the
input, an input line, and a status bar (FPS/frame time, line counts, active
filters, scroll position, contextual hint). No full-screen shade, drop shadow,
pulsing accent, or panel easing remain; only the open/close slide animates.
Output entries record a \texttt{ConsoleLevel}, a timestamp, and a repeat count
(identical consecutive lines collapse). The buffer holds 2,048 entries;
command recall keeps 256. New built-ins: \texttt{filter}, \texttt{loglevel},
\texttt{timestamps}, \texttt{copy}, \texttt{savelog}, \texttt{exec},
\texttt{bind}/\texttt{unbind} (F1--F12 except host-reserved F3, F5, F6, F11),
\texttt{!!}/\texttt{!n}/\texttt{!prefix} recall, and \texttt{help <word>}
search; a follow-up added Ctrl+R reverse search, \texttt{add},
\texttt{cycle}, \texttt{watch}/\texttt{unwatch} (a live strip of up to eight
variables), \texttt{writeconfig}, click-to-recall on echoed commands,
severity marks, filter-match highlights, and a closed-console error/warning
badge (\texttt{alerts}). Clipboard and file access go through
\symbolref{CommandLineCore} virtual hooks and are compiled out of WASM
guests.\\
\textbf{Why:} The console is a debugging tool; styling it like the product
menus blurred that line, and a debugging session needs severity filtering,
repeat collapsing, scripts, and log export more than decoration.\\
\textbf{Consequences:} Hidden entries wrap to zero lines, so D-UI2 scroll
metrics and D-P2 settled replay are unchanged; a scrolled-back view stays
anchored while output arrives. Log prefixes (\texttt{ERROR:} and so on) stay
in entry text for compatibility. Guest console bridges forward collapsed
repeats and restore the success level on the host.\\
\textbf{Supersedes:} the console part of D-R20 (UiTheme styling)
\end{decisionbox}

%----------------------------------------------------------------------------
\subsection*{D-UI5 --- Detached console window without a second GL context}
\begin{decisionbox}[title={D-UI5 \quad 2026-09-23}]
\textbf{Decision:} The developer console can pop out into its own OS window
(header \texttt{pop out} button, \texttt{console\_mode detached}, or dragging
the floating title bar past the game window's edge) and return (\texttt{dock}
button, closing the window, or the console key from either window). The
window is
a GLFW \texttt{GLFW\_NO\_API} window owned by \symbolref{DebugModule}
(\texttt{Platform/PixelWindow}); \texttt{SoftwareCanvas} rasterizes the same
\symbolref{GameVisual} primitives the in-game console submits, and the
platform presents the CPU image (Windows GDI; unsupported elsewhere).
\symbolref{CommandLine} only raises detach/dock/close requests and composes
at the detached window's size; while detached \texttt{isOpen} is false so the
game keeps its input.\\
\textbf{Why:} A second OpenGL context would break the single-context backend
contract (unshared VAOs, cached device state); a CPU presenter keeps the
renderer untouched and the console's look identical.\\
\textbf{Consequences:} Console composition is split from token submission.
Detached drawing costs a CPU raster only when composition is dirty (typing,
caret blink, output). Linux reports detaching as unavailable. Headless tests
cover the canvas and request logic; the OS window itself needs a manual
Windows smoke.
\end{decisionbox}

%----------------------------------------------------------------------------
\subsection*{D-R11 --- Backend constructed at composition root}
\begin{decisionbox}[title={D-R11 \quad 2026-08-06}]
\textbf{Decision:} Construct the concrete production backend
(\symbolref{GLBackend}) through fallible \symbolref{Illumo::initialize} and
inject it into \symbolref{Renderer} as \texttt{std::unique\_ptr<IBackend>}.
The factory constructs only; the host owns the backend's single initialization
attempt and reports failure without terminating the process. \symbolref{Renderer}
must not include or construct OpenGL types. Tests inject
\symbolref{MockBackend} through the same ownership path.\\
\textbf{Why:} Dependency inversion: \texttt{Renderer $\rightarrow$ IBackend},
composition root $\rightarrow$ concrete backend. Production and headless test
architectures become structurally identical.\\
\textbf{Consequences:} \pathref{Renderer.h} no longer includes
\pathref{OpenGL/GLBackend.h}; OpenGL coupling stays under
\pathref{Rendering/OpenGL/} and the host composition file.
\end{decisionbox}

%----------------------------------------------------------------------------
\subsection*{D-R12 --- CommandQueue overflow policy}
\begin{decisionbox}[title={D-R12 \quad 2026-08-06}]
\textbf{Decision:} Keep the fixed 2048-token \symbolref{CommandQueue} capacity.
When full, drop additional tokens, count drops, and log an error once per
frame (once per \texttt{Reset} interval) rather than failing silently or
growing the buffer unbounded.\\
\textbf{Why:} Silent loss hid runaway emitters; unbounded growth is not
acceptable on the production frame path. One log per overflowing frame is
visible without spamming every discarded command.\\
\textbf{Consequences:} Callers can query drop counters in tests; capacity
changes require profiling evidence, not opportunistic growth.
\textbf{Status:} Fixed-capacity policy superseded by bounded vector growth in
D-R16; one-warning and rejection-metric intent remains.
\end{decisionbox}

%----------------------------------------------------------------------------
\subsection*{D-WW1 --- Wireworld seed and edit brush}
\begin{decisionbox}[title={D-WW1 \quad 2026-08-06}]
\textbf{Decision:} Startup seeding is ruleset-aware. Binary life-like modes seed
a centered Game-of-Life glider. Wireworld seeds a horizontal conductor line
with a head+tail electron pair. In Wireworld edit mode, left-paint uses a
sticky brush selected by keys \texttt{1}/\texttt{H} (head), \texttt{2} (empty),
\texttt{3}/\texttt{T} (tail), \texttt{4} (conductor); right-click paints empty.\\
\textbf{Why:} The previous GoL-only glider produced five heads under Wireworld,
and mouse paint could not place heads without the console \texttt{setcell}
command.\\
\textbf{Consequences:} Wireworld is usable for circuit editing without leaving
edit mode; focused tests cover seed cells and brush selection state.
\end{decisionbox}

%----------------------------------------------------------------------------
\subsection*{D-C2 --- CellGrid domain extracted from Canvas}
\begin{decisionbox}[title={D-C2 \quad 2026-08-06}]
\textbf{Decision:} Introduce \symbolref{CellGrid} as pure dense simulation
storage (life buffer + dirty tracking) with no \symbolref{Renderer}, window, or
camera dependency. \symbolref{Canvas} extends \symbolref{CellGrid} for
presentation (RGB fade, texture upload) and GPU enrollment.
\symbolref{RuleSet} operates on \texttt{CellGrid*} only.\\
\textbf{Why:} The architecture assessment required a truthful domain boundary so
simulation can be constructed, stepped, and read without a graphics stack.
D-C1's intentional monolith is refined, not discarded: presentation still lives
on \symbolref{Canvas}.\\
\textbf{Consequences:} Headless domain tests construct \symbolref{CellGrid}
(or a domain-only \symbolref{Canvas} with null render services) and drive
shipped \symbolref{RuleSet} implementations. A further
\texttt{CanvasRenderer} split remains optional.\\
\textbf{Supersedes / refines:} D-C1 (monolith until pain).
\end{decisionbox}

%----------------------------------------------------------------------------
\subsection*{D-C3 --- Sparse infinite production canvas}
\begin{decisionbox}[title={D-C3 \quad 2026-08-07}]
\textbf{Decision:} Replace the finite dense production path with
\symbolref{SparseCellGrid}: signed 64-bit world coordinates, non-background
16$\times$16 chunks, temporary 18$\times$18 read-only halos, and serial
non-toroidal stepping. Present it through a bounded \symbolref{CanvasView}
with one reusable RGB staging texture and one quad. Always write
version 2 sparse saves; read both version 2 and the prior dense format after
validating temporary state.\\
\textbf{Why:} The finite canvas was the remaining product-scale boundary:
editing, simulation, camera, rendering, and persistence all needed one
coherent world-coordinate model rather than a dormant capped prototype.\\
\textbf{Consequences:} Empty space is implicit, negative coordinates are
first-class, far-apart chunks are not capped, and newly revealed view cells
snap before visible-area fading. Dense \texttt{CellGrid}/\texttt{Canvas} code
remains only for compatibility tests and legacy import; it is not a second
runtime path.\\
\textbf{Supersedes:} D-C1 and D-C2 for the production game path; their dense
domain remains compatibility material.
\end{decisionbox}

%----------------------------------------------------------------------------
\subsection*{D-C4 --- Cell-aligned sparse presentation}
\begin{decisionbox}[title={D-C4 \quad 2026-08-07}]
\textbf{Decision:} Render \symbolref{CanvasView} as a world-space quad whose
nearest-filtered texture texels each represent one 16$\times$16 world cell.
Grow staging capacity as needed to cover the camera viewport. The editor cursor
uses the same centered cell bounds.\\
\textbf{Why:} A screen-space, linearly filtered texture decoupled display
texels from world cells, blurring colors and misaligning the cursor after pan
or zoom.\\
\textbf{Consequences:} Cell edges remain sharp and cursor, paint, and display
refer to the same cell. Zoomed-out views retain a bounded texture without
using linear filtering to hide undersampling.\\
\textbf{Refines:} D-C3 presentation details.
\end{decisionbox}

\subsection*{D-R13 --- Single production frame path}
\begin{decisionbox}[title={D-R13 \quad 2026-08-06}]
\textbf{Decision:} \symbolref{Illumo::render} always rebuilds the per-frame
drawable list and submits tokens through \symbolref{Renderer::RenderScene}.
Remove the env \texttt{UseTokenProof} product bypass. Keep
\symbolref{Renderer::RenderProofQuad} for headless token e2e tests only.
Hybrid immediate \texttt{Draw()} remains only as a stub/test fallback when
\texttt{AppendCommands} returns false.\\
\textbf{Why:} Competing production frame paths create ambiguous ordering and
test gaps.\\
\textbf{Consequences:} One default product path; proof coverage stays in
\texttt{IllumoTests} without affecting the live loop.
\end{decisionbox}

%----------------------------------------------------------------------------
\subsection*{D-OWN1 --- Module-owned mode splash}
\begin{decisionbox}[title={D-OWN1 \quad 2026-08-06}]
\textbf{Decision:} Own the EDIT/NORMAL splash label as
\texttt{std::unique\_ptr<SplashText> modeSplash} on
\symbolref{CellGameModule}. Delete the file-scope \texttt{stateSplash}
mutable global.\\
\textbf{Why:} File-scope mutable product UI hid ownership and complicated
teardown and multi-instance reasoning.\\
\textbf{Consequences:} Splash lifetime follows the game module; platform
\texttt{SaveLoad} dialog entry points remain platform implementations without
becoming a new engine singleton.
\end{decisionbox}

%----------------------------------------------------------------------------
\subsection*{D-C5 --- Revision-gated adaptive sparse overview}
\begin{decisionbox}[title={D-C5 \quad 2026-08-07}]
\textbf{Decision:} Keep nearest-filtered, one-texel-per-cell presentation while
the visible world fits the active texture. At far zoom, cap the active texture
to the larger of configured canvas dimensions or roughly one texel per four
screen pixels, and accumulate source-cell palette density into each overview
texel. Rebuild only on a sparse-grid revision, palette update, or visible-region
change; traverse only intersecting sparse chunks.\\
\textbf{Why:} Expanding the texture and hash-looking up every visible cell on
every render frame made wide camera views lag even when the displayed world was
unchanged. Linear filtering would hide neither the CPU work nor the cell model.\\
\textbf{Consequences:} Near-view cell edges, cursor alignment, and paint stay
exact. Far views remain crisp rather than blurred, retain sparse activity and
multi-state proportions, and use bounded fade/upload work. The texture budget
does not cap chunks, simulation cells, randomization source bounds, or saves.\\
\textbf{Refines:} D-C4 presentation scaling details.
\end{decisionbox}

%----------------------------------------------------------------------------
\subsection*{D-P8 --- Responsive parallel sparse stepping}
\begin{decisionbox}[title={D-P8 \quad 2026-08-07}]
\textbf{Decision:} Keep sparse target discovery and next-map merge serial and
deterministic. For 32 or more target chunks, \symbolref{SparseCellGrid} uses a
grid-owned reusable pool of up to eight workers to evaluate independent
18$\times$18 halo / 16$\times$16 core results into separate slots. Normal mode
performs at most two generations per rendered frame and drops further catch-up
debt.\\
\textbf{Why:} Dense random sparse worlds are generation-bound after the direct
halo lookup optimization. Per-generation thread creation would erase much of
the gain, while unrestricted catch-up makes camera and input visibly stutter.\\
\textbf{Consequences:} Rule evaluation, persistence ordering, revisions, and
the bounded presentation path remain deterministic. Small patterns remain
serial. \texttt{SparseAdvanceStats}, Tracy zones, serial/parallel identity
tests, and \texttt{IllumoGame.Sim.SparseMicroBench} make the threshold observable.
\end{decisionbox}

%----------------------------------------------------------------------------
\subsection*{D-P9 --- Adaptive cell candidates for sparse chunks}
\begin{decisionbox}[title={D-P9 \quad 2026-08-08}]
\textbf{Decision:} Keep 16$\times$16 chunks as authoritative storage, but
maintain a four-word occupancy mask per stored chunk. Below an average of 48
occupied cells per active chunk, evaluate only each non-background cell plus
the eight neighbors of counted cells, accumulating neighbor counts in temporary
candidate records. At or above that density, retain D-P8's direct 18$\times$18
halo evaluator and bounded worker pool.\\
\textbf{Why:} A wide thin colony made full 16$\times$16 core evaluation spend
most of its time on unchanged background cells. Sparse storage alone did not
reduce that per-chunk work.\\
\textbf{Consequences:} Every shipped ruleset has a stable empty background, so
untouched cells remain implicit. Candidate results are byte-identical to the
full-chunk evaluator; tests cover wide colonies, multi-state transitions, and
dense fallback. Status and Tracy distinguish active cells, candidates, and the
selected evaluation path.
\end{decisionbox}

%----------------------------------------------------------------------------
\subsection*{D-P10 --- Retained chunk-local candidate scratch}
\begin{decisionbox}[title={D-P10 \quad 2026-08-08}]
\textbf{Decision:} Replace the per-world-cell candidate hash map and diagnostic
candidate-chunk hash set with two grid-owned, capacity-retaining vectors. Sort
the exact affected chunk addresses, then use one contiguous scratch record per
target chunk containing a 256-bit candidate mask and 256 neighbor counts. Reset
record contents with \texttt{clear}/\texttt{resize} while retaining the vectors'
high-water capacity.\\
\textbf{Why:} Profiling a wide sparse colony found hundreds of thousands of
temporary hash-node allocations per generation; hash construction and teardown
consumed substantially more time than rule evaluation.\\
\textbf{Consequences:} Candidate discovery performs no per-world-cell hash
allocation, scratch records are contiguous, and unchanged colony width reuses
capacity. The authoritative sparse chunk map and full-chunk fallback remain
unchanged. Tests cover exact candidate counts, capacity reuse, boundary behavior,
and byte identity across all shipped rulesets.\\
\textbf{Refines:} D-P9 candidate storage.
\end{decisionbox}

%----------------------------------------------------------------------------
\subsection*{D-P11 --- Generation-stamped flat candidate index}
\begin{decisionbox}[title={D-P11 \quad 2026-08-08}]
\textbf{Decision:} Retain a power-of-two, open-addressed table from signed
chunk address to candidate-scratch index. Use linear probing and a 64-bit
generation stamp so prior slots become logically empty without clearing the
table. Create affected chunks directly during collection, retain the contiguous
scratch vector, and grow both structures only beyond their previous high-water
marks.\\
\textbf{Why:} After D-P10, the 16,384-colony benchmark still performed 294,912
binary searches over 65,536 large scratch records per generation, while address
collection/sorting and scratch construction dominated measured simulation time.\\
\textbf{Consequences:} Sparse candidate discovery no longer sorts, uniques, or
binary-searches chunk addresses. Signed unbounded coordinates remain intact;
capacity is reused across generations; and focused tests cover table growth,
generation reset, negative coordinates, and byte identity against full-chunk
evaluation for every shipped ruleset. The 16,384-colony Release benchmark
improved from roughly 18 to roughly 40 generations per second on the measured
machine.\\
\textbf{Refines:} D-P10 candidate-address discovery; the candidate masks and
neighbor counts remain chunk-local and contiguous.
\end{decisionbox}

%----------------------------------------------------------------------------
\subsection*{D-P12 --- Retained transactional chunk-map nodes}
\begin{decisionbox}[title={D-P12 \quad 2026-08-08}]
\textbf{Decision:} Keep the authoritative sparse chunk map and a grid-owned
inactive next-generation map. Before each build, extract inactive map nodes into
a retained vector of node handles. Reassign and reinsert those nodes for output
chunks, retaining both maps' bucket capacity. Allocate a new node only when the
output exceeds the previous node high-water count. Use the same lifecycle for
cell-candidate and full-chunk result paths.\\
\textbf{Why:} After candidate scratch became allocation-free, the 16,384-colony
case still allocated and freed about 32,775 nodes and moved roughly 11 MB through
the heap per generation. Most were authoritative output-map nodes.\\
\textbf{Consequences:} A steady 16,384-colony generation now reports zero output
chunk-node allocations while reusing and retaining 32,768 nodes. The retained
and allocating A/B benchmark keeps identical rules and map semantics and shows
the retained path consistently faster on the measured machine. The current map
remains untouched until result evaluation succeeds and content comparison chooses
the transactional swap; unchanged generations keep their revision. Status and
Tracy expose allocation, reuse, retention, preparation, recycling, and finish
work.\\
\textbf{Refines:} D-P10 and D-P11 generation-local storage; authoritative sparse
map semantics remain unchanged.
\end{decisionbox}

%----------------------------------------------------------------------------
\subsection*{D-P13 --- Coarse parallel candidate evaluation}
\begin{decisionbox}[title={D-P13 \quad 2026-08-08}]
\textbf{Decision:} Once serial candidate scratch construction completes, divide
large candidate sets into retained contiguous ranges containing roughly 2,048
candidate cells. Use the existing reusable pool to evaluate ranges into
independent retained result slots, then merge results serially through D-P12's
transactional node recycler. Enable automatic parallelism at 16,384 candidate
cells and retain direct serial evaluation below that threshold.\\
\textbf{Why:} Candidate rule evaluation remained serial after flat-index and
node-allocation work. Claiming one atomic job per mostly empty target chunk would
make 65,536 claims for the 16,384-colony benchmark, so work is balanced by actual
candidate-cell count instead.\\
\textbf{Consequences:} The 16,384-colony case forms about 120 jobs rather than
65,536. Explicit four- and eight-worker Release benchmarks improve over retained
serial evaluation on the measured machine; automatic execution is capped at
four workers because additional workers gave little benefit on this path. A
256-colony case keeps its original direct loop and avoids result-buffer overhead.
Tests verify coarse job granularity, requested worker use, and byte identity with
serial candidates. Status and Tracy report range count, workers, parallel
evaluation, and serial merge work.\\
\textbf{Refines:} D-P9 candidate evaluation and D-P8 worker-pool use.
\end{decisionbox}

%----------------------------------------------------------------------------
\subsection*{D-P14 --- Retained changed-region frontier}
\begin{decisionbox}[title={D-P14 \quad 2026-08-08}]
\textbf{Decision:} Preserve the inactive chunk map as the complete prior
generation and retain generation-stamped flat sets for changed chunks and their
next targets. Expand each changed chunk by its eight neighbors. An empty set
returns without evaluating cells; at most 64 targets are evaluated and patched
transactionally into the prior map. Cell edits and ruleset-type changes mark the
affected chunks. More than 64 changed chunks invalidates the optimization for
that generation and selects the existing complete candidate or halo path.\\
\textbf{Why:} For deterministic radius-one rules, a chunk outside the changed
set and its neighbors sees the same input neighborhood as the prior generation
and therefore cannot change. Complete candidate discovery still rescanned every
stable chunk in a wide mostly-static world.\\
\textbf{Consequences:} Settled still lifes advance with zero target evaluation.
The measured 1,024-block world with one blinker evaluates 15 of 1,026 active
chunks and runs roughly 15 times faster than complete candidate discovery.
Broad moving colonies retain the bounded tracking cost and fall back to the
already parallel complete paths. Tests cover settled stepping, edits, ruleset
changes, and byte identity against forced full evaluation.\\
\textbf{Refines:} D-P12's retained prior map and the D-P9/D-P13 complete sparse
paths.
\end{decisionbox}

%----------------------------------------------------------------------------
\subsection*{D-P15 --- Per-target counted-state adaptation}
\begin{decisionbox}[title={D-P15 \quad 2026-08-08}]
\textbf{Decision:} Store separate 256-bit masks for non-background cells and
cells that contribute to Moore-neighbor counts. When candidate scratch is
built, classify every target independently: use its exact candidate mask below
the counted-neighbor contribution threshold, otherwise use the direct
18$\times$18 halo evaluator. Work ranges charge halo targets as 256 cells.
Worlds whose source chunks are all counting-dense bypass scratch construction.\\
\textbf{Why:} A global average misclassified mixed dense/sparse worlds. It also
treated every stored Wireworld conductor as neighbor-counting work even though
only electron heads contribute to conductor transitions.\\
\textbf{Consequences:} One generation may contain both candidate and halo
targets while preserving deterministic transactional output. The measured
9$\times$9 all-conductor Wireworld case contains 20,736 stored cells but zero
counted cells; candidate stepping measured roughly 8,019 generations per second
versus 4,249 for forced halos on the development machine. Tests compare mixed
per-target output and conductor-only output against forced complete halos.
Status exposes stored/counted cells, candidate/halo target totals, and the mixed
path.\\
\textbf{Refines:} D-P9's global density selection and D-P13's candidate work
ranges.
\end{decisionbox}

%----------------------------------------------------------------------------
\subsection*{D-P16 --- Cached transitions and rolling halo stencil}
\begin{decisionbox}[title={D-P16 \quad 2026-08-09}]
\textbf{Decision:} Lazily build an immutable 256-state by 9-neighbor transition
table for each ruleset instance before worker dispatch. Candidate, halo, and
compatibility dense loops index that table. Reduce every 18$\times$18 halo with
three rolling horizontal count rows; each output cell sums those rows and
removes the center state instead of rescanning eight neighbors.\\
\textbf{Why:} The transition domain is only 2,304 inputs, while virtual calls,
rule branches, and repeated Moore-neighborhood comparisons occurred for every
evaluated cell in every generation. The table preserves the existing ruleset
interface while moving that work out of the hot loops.\\
\textbf{Consequences:} A ruleset performs exactly 2,304 virtual transitions on
first use and none during subsequent generation evaluation. On the development
machine, the 24$\times$24-chunk benchmark improved from roughly 419 to 673
serial generations per second and from 1,146 to 1,605 with four workers. Tests
compare every table entry for all seven shipped rules and keep candidate, halo,
serial, and parallel outputs byte-identical. No SIMD or GPU simulation is added.\\
\textbf{Refines:} D-P8 dense halo evaluation and all D-P9--D-P15 sparse paths.
\end{decisionbox}

%----------------------------------------------------------------------------
\subsection*{D-P17 --- Transactional sparse population aggregates}
\begin{decisionbox}[title={D-P17 \quad 2026-08-09}]
\textbf{Decision:} Cache stored-cell and counted-cell totals per chunk and
maintain aggregate stored, counted, and candidate-preferred chunk totals for
both the authoritative and retained inactive maps. Update the aggregates in
the same operations that edit, assign, clear, patch, insert, recycle, or swap
chunk data. Read the authoritative aggregates at the start of
\symbolref{SparseCellGrid::advance} instead of popcounting every chunk.\\
\textbf{Why:} The retained changed frontier already made settled worlds perform
no cell evaluation, but \texttt{advance} still scanned and popcounted every
allocated chunk before reaching that return. Large static worlds therefore
paid work proportional to their stored chunk count on every requested step.\\
\textbf{Consequences:} Settled stepping and complete-path selection are
constant-time in allocated chunk count before any actual target construction.
Status retains exact stored/counting totals and now exposes the cached
candidate-preferred chunk count. Focused tests cover direct state changes,
threshold crossings, chunk replacement, clear, full-map and frontier-map
publication, and grid swaps. Each stored chunk adds two 16-bit counters and
each retained map adds three aggregate counters.\\
\textbf{Refines:} D-P14's empty-frontier fast path and D-P15's separate stored
and counted masks.
\end{decisionbox}

%----------------------------------------------------------------------------
\subsection*{D-P18 --- Cost-adaptive candidate frontiers}
\begin{decisionbox}[title={D-P18 \quad 2026-08-09}]
\textbf{Decision:} Replace D-P14's 64-target selection cutoff with a work-unit
comparison. Retain up to 16,384 changed chunk addresses, expand them by one
chunk, and build exact candidate masks from only the sparse source chunks that
can affect those targets. Classify frontier targets as candidates or halos by
counted-neighbor contributions. Compare frontier target/source bookkeeping,
candidate construction, contributions, and evaluation cells with a
complete-path estimate derived from D-P17's cached totals; patch the retained
map when frontier work is no greater.\\
\textbf{Why:} Eight separated localized changes could expand beyond 64 targets
and fall abruptly into complete-world preparation even when most of a large
world was static. The target count also treated a sparse candidate target as if
it required all 256 halo cells. The earlier 4,096-address journal was then
cleared for the rest of a dense generation, so later shrinking activity could
not return to frontier tracking without rebuilding from a complete scan.\\
\textbf{Consequences:} Local work can cross the former cutoff without a path
cliff, while broad separated changes are cost-rejected. Sparse frontier
evaluation shares candidate/halo semantics and worker ranges with complete
mixed stepping. Tracking and evaluation are separate: the 16,384-address cap
bounds persistent bookkeeping, while the work comparison still selects frontier
versus complete evaluation. Overflowing that cap or failing allocation
invalidates the journal; a later smaller generation can use frontier again.
Focused coverage exercises more than 64 frontier targets, candidate selection,
forced-complete byte identity, local-source bounds, broad-change fallback, a
burst above 4,096 changed chunks, and return to an empty frontier.\\
\textbf{Refines:} D-P14 changed-region selection and D-P15 per-target
adaptation.
\end{decisionbox}

%----------------------------------------------------------------------------
\subsection*{D-P19 --- Parallel candidate preparation without blanket clears}
\begin{decisionbox}[title={D-P19 \quad 2026-08-09}]
\textbf{Decision:} Derive candidate target edges and corners directly from the
four 64-bit counting-mask words. Do not value-initialize all 256 neighbor
counters in each scratch record; initialize a counter when its candidate bit
is first enrolled. Keep the locality-friendly source-centric builder for
serial preparation. For large automatic or explicitly parallel preparation,
assign independent target scratch records to the retained worker pool; each
worker reads the target's 3-by-3 source neighborhood and writes only its own
record. Report preparation workers separately from evaluation workers.\\
\textbf{Why:} Candidate collection previously rescanned every counted cell to
discover boundary participation, cleared 256 bytes per target even though only
candidate entries are read, and finished all scratch construction before
workers could begin. A fully target-centric serial traversal regressed sparse
locality, so preparation needs distinct serial and parallel layouts.\\
\textbf{Consequences:} Small sparse worlds retain the faster serial traversal
while avoiding blanket counter stores. Large worlds can parallelize candidate
construction without shared destination writes, using the same retained pool
as evaluation. Focused tests cover signed edge/corner targets, forced one- and
four-worker preparation, scratch reuse, and byte identity against complete
halos.\\
\textbf{Refines:} D-P12 retained candidate scratch and D-P13 coarse candidate
worker ranges.
\end{decisionbox}

%----------------------------------------------------------------------------
\subsection*{D-P20 --- Retained complete halos from counting masks}
\begin{decisionbox}[title={D-P20 \quad 2026-08-09}]
\textbf{Decision:} Replace the complete-halo path's generation-local target
hash set, address/result vectors, and target sort with a retained
generation-stamped flat target index plus retained target and result vectors.
For evaluation, extract one 18-bit counted row from the 3-by-3 source chunks'
counting masks and feed three rolling horizontal rows directly into the
transition loop. Read cell-state bytes only from the target core; do not
materialize a temporary 18-by-18 byte halo.\\
\textbf{Why:} Dense fallback paid repeated allocation, hashing, and sorting
before evaluation, then populated 324 byte states per target even though the
neighbor loop only needed counted-state membership around the core. This work
was visible both for all-dense worlds and for halo targets selected inside
mixed generations.\\
\textbf{Consequences:} Complete target/index/result capacity survives at its
high-water mark, while generation stamps make prior target slots logically
empty. Halo evaluation keeps the existing serial/worker split and exact
transition semantics but removes byte-halo construction. The informational
bench improves both serial and four-worker dense stepping; focused coverage
checks capacity reuse, index generations, signed boundaries, multistate rules,
and candidate/full-halo byte identity.\\
\textbf{Refines:} D-P8 reusable halo workers and D-P16 rolling halo counts.
\end{decisionbox}

%----------------------------------------------------------------------------
\subsection*{D-P21 --- Changed-tile sampling and active fades}
\begin{decisionbox}[title={D-P21 \quad 2026-08-09}]
\textbf{Decision:} Publish a revision-scoped retained list of current or
removed chunks after a single edit or completed generation. When
\symbolref{CanvasView} displays one texel per cell and exactly one revision is
unseen, resample only the visible portions of those 16-by-16 tiles. Retain a
deduplicated set of texels whose display RGB differs from target RGB; fade
ticks and zero-speed snaps visit only that set. Treat repeated unchanged
\texttt{setFadeSpeed(0)} calls as no-ops.\\
\textbf{Why:} Any content revision previously rebuilt the whole active view,
and any active fade scanned every visible texel each frame. The normal module
also reapplied the environment fade value every frame, causing a full snap
scan whenever that value was zero.\\
\textbf{Consequences:} Common near-zoom edits and generations perform work per
visible changed chunk, while fade work follows colors still moving. Density
overviews snap instead of fading. Bins that cover at least a quarter of the
cache stop further bin marking and resample the complete bounded cache.
Revision gaps, camera or palette changes, whole-grid replacement, and
invalid change journals use the complete bounded resample.
The status command reports fading, last-sampled, and last-fade-visited texels.
Focused tests cover edits, removals, generations, replacement invalidation,
overview snap, dense complete-sample fallback, revision-gap fallback, active
fade counts, and repeated zero-speed
configuration.\\
\textbf{Refines:} D-P6 CPU RGB fade and D-P14 retained changed chunks.
\end{decisionbox}

%----------------------------------------------------------------------------
\subsection*{D-P22 --- Measured synchronous simulation budget}
\begin{decisionbox}[title={D-P22 \quad 2026-08-09}]
\textbf{Decision:} Keep one due generation guaranteed in NORMAL mode. Before
starting a second synchronous generation, measure the first and require the
elapsed frame simulation time plus that observed step cost to fit inside a
4\,ms simulation slice. Retain the two-step ceiling and fractional accumulator
remainder while dropping further catch-up debt. Track a smoothed achieved TPS
separately from requested \texttt{tps} $\times$ \texttt{speedFactor}, plus the
longest step and total simulation time in the current frame.\\
\textbf{Why:} A count-only two-step cap still allowed two expensive generations
to execute before rendering. The configured TPS described demand, not the rate
the main thread actually sustained, so status could not distinguish scheduling
pressure from simulation cost.\\
\textbf{Consequences:} Cheap worlds can still take the bounded second step,
while a costly first generation yields to input, presentation, and rendering.
A single costly generation remains synchronous and visible in timing
diagnostics; background snapshot publication is deferred unless profiling shows
its p95 still exceeds the render-frame budget after CPU-path improvements.
Focused coverage forces zero and ample budget values to verify the one-step
guarantee, second-step gate, debt policy, and status diagnostics.\\
\textbf{Refines:} D-P8 responsive parallel sparse stepping.
\end{decisionbox}

%----------------------------------------------------------------------------
\subsection*{D-P23 --- Retained texture capacity and explicit release}
\begin{decisionbox}[title={D-P23 \quad 2026-08-09}]
\textbf{Decision:} When one CanvasView texture dimension grows by a small
amount, allocate the greater of the required size or 150\% of its current
capacity. Keep the opaque texture handle stable across re-enrollment, but make
\symbolref{GLBackend} explicitly destroy the old registry entry before
installing its replacement. Make \symbolref{GLTexture} destruction idempotently
delete its GL texture and two PBOs. Add \texttt{DestroyTexture} to
\symbolref{IBackend}/\symbolref{Renderer} and release the display texture from
\symbolref{CanvasView}'s destructor.\\
\textbf{Why:} Smooth zoom could increase the active dimensions by a few texels
on consecutive frames, repeatedly reallocating four CPU buffers and
re-enrolling the GPU texture. Registry replacement destroyed a C++
\texttt{unique\_ptr}, but the empty GLTexture destructor leaked the old texture
and any PBO pair. The final live texture also remained registered until global
backend shutdown after its view was gone.\\
\textbf{Consequences:} Nearby zoom sizes reuse retained CPU/GPU capacity, while
large jumps allocate their exact required bound. Replacement and view teardown
now reclaim GL resources with the context-owning backend. Shutdown remains an
idempotent bulk fallback; general refcounts and generational handle reuse are
still deferred. Headless coverage verifies retained capacity, stable handles,
and explicit view release; the full application build compiles the OpenGL
replacement path.\\
\textbf{Refines:} D-R4 shutdown lifecycle and D-C5 bounded overview capacity.
\end{decisionbox}

%----------------------------------------------------------------------------
\subsection*{D-P24 --- Exact padded camera cache and stable LOD}
\begin{decisionbox}[title={D-P24 \quad 2026-08-09}]
\textbf{Decision:} Separate visible-resolution diagnostics from a globally
aligned sampled cache padded by two 16-cell chunks on every side. Keep exact
one-texel-per-cell output near zoom. Use integer density LOD far away, coarsen
immediately to fit, and refine only when the next level fits within 80\% of the
budget. Map one-revision changed chunks to deduplicated exact or overview bins.
When LOD, size, and alignment are unchanged, shift the cache origin by copying
retained CPU texels and resampling only newly exposed strips.\\
\textbf{Why:} Continuous pan and smooth zoom previously changed source bounds
often enough to resample and reapply the complete active view every frame.
Following a compact moving pattern still crossed the two-chunk pad every few
frames and paid a full refill even though most texels were already sampled.\\
\textbf{Consequences:} Sub-cache camera motion changes only the MVP. Aligned
origin shifts keep the overlap and sample strips; LOD/resize, palette changes,
revision gaps, torus wrap, replacement, and non-aligned jumps still trigger
exact bounded refills. Density overviews snap color changes instead of
enrolling fade texels, and dense one-revision bins that cover a quarter of the
cache resample the complete cache. Status exposes visible/cache dimensions,
LOD, refill count, and rolling refill latency; deterministic and warmed camera
tests cover reuse, scroll overlap, boundaries, hysteresis, overview snap, and
overview updates.\\
\textbf{Refines:} D-C5 bounded overview and D-P21 changed-tile sampling.
\end{decisionbox}

%----------------------------------------------------------------------------
\subsection*{D-P25 --- Bounded tiled uploads and non-waiting PBO ring}
\begin{decisionbox}[title={D-P25 \quad 2026-08-09}]
\textbf{Decision:} Track dirty 16-by-16 texel tiles, merge adjacent runs into at
most eight rectangles when they cover at most half their enclosing AABB, and
otherwise submit the AABB. Upload rectangles through 64\,KiB directly. For
larger requests, select a signaled slot from three PBOs with \texttt{GLsync}
fences without waiting; use direct upload when every slot is busy or mapping
fails.\\
\textbf{Why:} One dirty AABB could copy a large unchanged gap, while mapping the
same PBO path synchronously could couple upload availability to frame latency.\\
\textbf{Consequences:} \texttt{RenderCommand::UpdateTexture} and the
renderer/backend boundary remain unchanged and OpenGL 3.3 is sufficient.
Texture replacement and destruction delete all fences and buffers. Tracy zones
separate selection, map, copy, unmap, and submission; status retains rolling
requested bytes and rectangle counts.\\
\textbf{Refines:} D-P4/D-P6 dirty upload staging and D-P23 explicit release.
\end{decisionbox}

%----------------------------------------------------------------------------
\subsection*{D-P26 --- Evidence-gated published simulation runner}
\begin{decisionbox}[title={D-P26 \quad 2026-08-09}]
\textbf{Decision:} Because the warmed 16,384-colony fixture exceeds the 4\,ms
p95 gate, keep two sparse grids and one persistent simulation runner. The main
thread owns the published grid while the worker exclusively mutates the spare.
Publish completed generations only at frame boundaries and carry a
\texttt{SparseGenerationDelta} so the former display grid can be synchronized
before reuse. Permit one request and no backlog; drop overdue whole steps while
retaining the fractional accumulator remainder.\\
\textbf{Why:} A single representative generation can exceed the complete
16.7\,ms frame budget even after CPU hot-path improvements, so any synchronous
policy still exposes that cost as input and render lag.\\
\textbf{Consequences:} Canvas cache reads never lock against simulation.
Achieved TPS counts completed publications. Pause/edit, persistence, ruleset
changes, manual stepping, and shutdown drain first. Snapshot records include
cell and occupancy/counting masks, broad mirrors update retained chunk nodes in
place, and the worker catches allocation failures without terminating the
process. Focused coverage checks deterministic two-generation alternation,
unchanged sparse chunks, debt dropping, state-transition drains, and shutdown
with work outstanding.\\
\textbf{Supersedes:} D-P22 measured second synchronous step.
\end{decisionbox}

%----------------------------------------------------------------------------
\subsection*{D-P27 --- Cell-precise frontier activity}
\begin{decisionbox}[title={D-P27 \quad 2026-08-11}]
\textbf{Decision:} Retain a 256-bit state-change mask and a separate
counting-status-change mask for every changed chunk. Always reevaluate the
changed chunk itself. Enroll a neighboring target only when a counting change
touches its shared edge or corner. Carry the masks through edits, complete and
frontier generations, sparse publication deltas, mirror application, and grid
swaps.\\
\textbf{Why:} Chunk-address-only tracking expanded every local change to a
3-by-3 target region even when activity stayed in the chunk interior. Stored
states such as Wireworld conductors also caused neighbor work despite not
contributing to Moore counts.\\
\textbf{Consequences:} An interior blinker in a large settled world evaluates
one target chunk; boundary activity enrolls only the touching targets, and
non-counting stored chunks do not create neighbor work. Status reports exact
state/counting changed-cell totals. Focused tests cover interior, edge, binary,
and conductor cases plus byte identity with complete stepping.\\
\textbf{Refines:} D-P14 retained frontiers and D-P15 separate counting masks.
\end{decisionbox}

%----------------------------------------------------------------------------
\subsection*{D-P28 --- Adaptive exact halo memoization}
\begin{decisionbox}[title={D-P28 \quad 2026-08-11}]
\textbf{Decision:} Memoize halo results on demand with the complete
18-by-18 byte input state as the key. Use nine thread-exclusive shards (main
plus eight worker slots), 512 four-way entries per shard, hash-based set
selection, and mandatory full-key comparison. Sample up to 16 targets per shard
and activate at a 25\% hit rate. After three active generations below 10\%,
cool down for 32 generations. Bypass fewer than 32 halo targets and every
candidate-only workload; invalidate on ruleset type or transition-table
revision changes.\\
\textbf{Why:} Repeated isolated chunks often present identical halo states,
making neighbor reduction and 256 transition evaluations redundant. A global
cache or hash-only acceptance would add contention or risk incorrect output.
Unique worlds must not pay permanent key-building overhead.\\
\textbf{Consequences:} Each grid allocates shards lazily and remains below the
dual-grid 8\,MiB budget. Exact cached and uncached generations are byte
identical, including after ruleset changes. Three Release samples of 256
repeated dense chunks improved by 10.9--12.5\%; the 24-by-24 random control
changed by +0.8--1.3\%, and the candidate-only control's median change stayed
within 1\% with zero probes. Status exposes hits, probes, entries, bytes, and
activation state.\\
\textbf{Refines:} D-P16 cached transitions and D-P20 retained halo evaluation.
\end{decisionbox}

%----------------------------------------------------------------------------
\subsection*{D-P29 --- Allocation-free broad mirror replacement}
\begin{decisionbox}[title={D-P29 \quad 2026-08-11}]
\textbf{Decision:} Do not copy the completed sparse delta into the worker
grid's inactive-catch-up state. For incremental publication, use that grid's
existing changed-address journal and a retained index of incoming addresses;
skip an old address when the incoming record already supplies its final state.
For a full sparse replacement, extract every destination-map node into the
grid's retained recycler, rewrite keys and chunk data from the replacement
records, reinsert them, and accumulate aggregate statistics during insertion.
Report mirror, simulation advance, delta capture, and total worker latency as
separate Tracy zones and rolling status metrics.\\
\textbf{Why:} The 16,384-colony fixture exceeds the bounded frontier journal
and therefore publishes full sparse replacements. Its mirror path allocated a
temporary address set, performed per-record lookup/update, scanned the old map
to erase absent chunks, and then scanned again for statistics. Incremental
grids also retained a complete duplicate delta and could replay a state that
the next incoming delta immediately replaced.\\
\textbf{Consequences:} Full replacement performs one extraction pass and one
insertion pass with retained nodes and no temporary hash set. Incremental
catch-up retains only existing grid journals and writes each superseded address
once. Deterministic tests cover overlapping and non-overlapping incremental
changes plus a tracking-overflow full replacement. In two comparable
Release runs, dispersed mirror p95 fell from 17.7\,ms to 5.0--5.1\,ms and total
async p95 fell from 49.4\,ms to 28.3--28.6\,ms; the remaining dominant stage is
the 21.3--21.4\,ms generation advance.\\
\textbf{Refines:} D-P12 retained chunk nodes and D-P26 dual-grid publication.
\end{decisionbox}

%----------------------------------------------------------------------------
\subsection*{D-P30 --- Coarse direct-source candidate preparation}
\begin{decisionbox}[title={D-P30 \quad 2026-08-11}]
\textbf{Decision:} Keep the source-centric serial builder. For large or forced
parallel candidate preparation, link each target to its existing 3-by-3 source
chunks during target discovery. Store those nine pointer representations in
the first bytes of the otherwise uninitialized neighbor-count scratch, copy
them into worker-local storage before counting begins, and then reuse the bytes
for lazily initialized counts. Schedule retained target ranges sized to about
eight ranges per worker and no more than 256 targets, rather than claiming one
atomic work item per target. Keep the automatic candidate-worker cap at four.\\
\textbf{Why:} The 16,384-colony complete candidate path created 65,536 target
chunks. Target-parallel preparation repeated up to nine hash-map lookups and
one atomic claim for every target, while the later candidate evaluation was
already comparatively small. Adding permanent source-pointer arrays would
inflate every retained scratch record, including the serial path.\\
\textbf{Consequences:} The parallel path removes repeated source lookups and
reduces 65,536 work claims to 256 without increasing steady scratch-record
size. In comparable Release samples, candidate preparation fell from about
7.33\,ms to 1.58--1.83\,ms and four-worker throughput rose from about 48.5 to
68.9--69.5 generations per second; retained serial throughput remained about
43 generations per second. Dispersed asynchronous advance p95 fell from about
19.82\,ms to 14.74--15.53\,ms and total worker-cycle p95 from about 26.28\,ms to
21.46--22.29\,ms. An eight-worker production cap, keyed in-place output merge, and
worker-materialized output chunks were measured and rejected because they
worsened frame latency or total throughput. Stage timings and preparation and
evaluation range counts remain informational diagnostics.\\
\textbf{Refines:} D-P19 target-owned parallel preparation and D-P13 coarse
candidate evaluation.
\end{decisionbox}

%----------------------------------------------------------------------------
\subsection*{D-P31 --- Duplicate-first discovery and direct sparse output}
\begin{decisionbox}[title={D-P31 \quad 2026-08-16}]
\textbf{Decision:} Probe the retained candidate index before calculating
capacity. Check the 75\% load boundary and grow only after a probe finds a new
target. Count non-background output in independent evaluation ranges for exact
diagnostics, but retain target-sized bucket headroom. Before recycling every
inactive node, reset its aggregate statistics once. Candidate output acquires
and rekeys a recycled node first, constructs sparse cells directly in its
mapped storage, and adds statistics after construction. Keep dense halo output
on its complete-array bulk-copy path.\\
\textbf{Why:} The 16,384-colony path made 131,072 enrollment attempts for
65,536 unique targets, so half of the old capacity checks could never insert.
Sparse merge also cleared a temporary 256-cell chunk and then copied the whole
chunk into a node. Exact output reservation, however, reduced bucket headroom
and increased collision cost; sparse bit-walking also regressed dense output.\\
\textbf{Consequences:} Three same-machine Release baselines had medians of
52.6 generations per second, 7.06\,ms discovery, and 5.92\,ms merge. The scoped
variant measured 59.4 generations per second, 6.19\,ms discovery, and 5.86\,ms
merge. A direct-versus-temporary-copy A/B measured 5.86 versus 6.33\,ms median
merge. Across three warmed frame samples, asynchronous worker-cycle and advance
p95 medians improved by about 4.7\% and 4.9\%; synchronous p95 varied upward by
3.8\%, within the 5\% noise/regression gate. Status reports enrollments, index
growth, produced chunks, and change/recycle/output substages. Exact output
reservation and dense direct construction are rejected.\\
\textbf{Refines:} D-P11 candidate indexing, D-P12 retained nodes, and D-P30
direct-source preparation.
\end{decisionbox}

%----------------------------------------------------------------------------
\subsection*{D-P32 --- Synchronized presentation cadence}
\begin{decisionbox}[title={D-P32 \quad 2026-08-16}]
\textbf{Decision:} Default the Windows OpenGL path to swap interval one. Read
the persisted \texttt{vsync} value at swap boundaries and change GLFW's swap
interval only when that value changes; retain explicit uncapped mode for
profiling. In Debug, report frame-paced swap completions separately from CPU
frame submissions and show paced FPS as off while uncapped.\\
\textbf{Why:} Camera-only Tracy runs with swap interval zero completed 99\% of
CPU frames within 0.61\,ms and never exceeded 3.18\,ms, yet the overlay claimed
about 4,500--5,000 FPS while visible camera motion juddered. The same optimized
build with swap interval one locked to the 144\,Hz display: frame-start p50,
p95, and p99 were 6.94, 7.00, and 7.09\,ms, and p99.9 was 7.38\,ms. Four longer
intervals coincided with automation screenshot capture rather than camera
cache or upload work.\\
\textbf{Consequences:} Default motion is paced by the display instead of an
unbounded submission loop. Users can persist \texttt{vsync=0} or toggle it in
Debug for throughput profiling. The existing \texttt{IBackend::getFPS}
swap-completion counter supplies paced cadence only when synchronization is
active; the overlay never labels uncapped submissions as presented frames.
Headless coverage verifies the default/fallback policy and both label forms.\\
\textbf{Refines:} D-R13's single frame path and D-P24 camera-cache presentation.
\end{decisionbox}

%----------------------------------------------------------------------------
\subsection*{D-P33 --- Direct broad generations and exact topology reuse}
\begin{decisionbox}[title={D-P33 \quad 2026-08-17}]
\textbf{Decision:} Publish broad generations as a lightweight replacement
marker without chunk snapshots. Capture one-revision journals of at least
2,048 presentation chunks with the same marker. Let the spare grid read the immutable published
grid through \texttt{advanceFrom} and update its own authoritative chunk nodes
in place. Retain each target's 3-by-3 source pointers and an index-aligned
destination pointer. Give each grid an object-local topology epoch that changes
only when chunk presence or counted edge/corner participation changes; reuse the
candidate target index and pointers only when both source identity and epoch
match.\
\textbf{Why:} Full replacement capture and mirror work duplicated the broad
generation, while every dispersed step rediscovered 65,536 targets and looked
up 32,768 output nodes. The two alternating grids revisit the same parity of
oscillator topology, so exact object-local reuse is available even though the
logical world topology changes every generation.\
\textbf{Consequences:} Broad mirror and capture p50/p95 are zero in the
16,384-colony fixture. Candidate topology was reused in 327 of 330 warmed and
measured generations; a representative asynchronous p95 improved from about
17.7--20.3\,ms after direct publication to 9.7--11.3\,ms with topology and
output-slot reuse, while remaining byte-identical. An edit that changes
topology invalidates reuse before any retained pointer is read. A 32,768-chunk
contiguous-storage prototype made its isolated copy kernel 1.24 times faster,
but its estimated whole-generation ceiling was only about 0.6\,ms, so the
authoritative hash-map refactor is rejected. A two/four/eight-worker A/B measured
11.97/9.67/10.62\,ms p95, retaining the four-worker automatic cap. A paced Tracy
pan/zoom run at 144\,Hz measured 6.94/7.00/8.10\,ms frame-start p50/p95/max in
the refill window and only six refills across 36,892 frames.\
\textbf{Refines:} D-P26 publication, D-P30 source-linked preparation, D-P31
direct sparse output, and D-P32 paced presentation.
\end{decisionbox}

%----------------------------------------------------------------------------
\subsection*{D-C6 --- Configurable infinite or finite toroidal world}
\begin{decisionbox}[title={D-C6 \quad 2026-08-17}]
\textbf{Decision:} Keep \texttt{0 x 0} as the existing infinite,
non-toroidal sparse topology. Treat positive world width and height as a
finite chunk rectangle whose opposite axes are adjacent; reject mixed
zero/positive dimensions. Canonicalize public reads and writes and wrapped
neighbor enrollment into that rectangle. Expose ruleset, chunk dimensions,
TPS, simulation speed, fade speed, VSync, and fullscreen through a Game-owned
F1 \symbolref{ConfigurationMenu} compiled in Release and Debug. A topology
change drains simulation and starts a fresh centered world. Write topology in
sparse save version 3 while continuing to read version 2 and legacy dense
files as infinite worlds.\\
\textbf{Why:} Release builds previously had no supported settings surface, and
finite world dimensions are meaningful only when their edge semantics are
explicit. A torus preserves deterministic Moore neighborhoods without adding
hard-boundary rule exceptions.\\
\textbf{Consequences:} The optimized infinite sparse paths remain unchanged;
finite worlds use canonical wrapped candidate enrollment for binary and
multi-state rules. The bounded view clips presentation to the canonical finite
rectangle, leaves outside camera space background-colored, and conservatively
refills after finite-world changes. The menu remains primitive-composed through
\symbolref{GameVisual}, not a retained widget tree, with high-contrast labels
and contextual help. The menu Exit action requests window closure through
\symbolref{IRenderWindow}; the App main loop then performs normal shutdown.
Transient scalar animation state drives the eased entrance, staggered rows,
gliding selection, and value pulse without introducing widgets or delaying
input.
Version 3 topology, version 2 compatibility, invalid metadata, F1 integration,
and horizontal/vertical/corner seams have focused headless tests.\\
\textbf{Refines:} D-C3 sparse production state and D-R20 product UI.
\end{decisionbox}

%----------------------------------------------------------------------------
\subsection*{D-DOC2 --- Technical documentation and operational guidance}
\begin{decisionbox}[title={D-DOC2 \quad 2026-08-09}]
\textbf{Decision:} Keep canonical technical documentation under
\pathref{docs/}: current Markdown, package maps, dated records, provenance,
LaTeX sources, and generated PDF output. Treat \pathref{latex/illumo.tex} as
the canonical prose-book entrypoint and \pathref{latex/architecture-map.tex}
as the current chart-only entrypoint; \pathref{docs/build.ps1} builds both.
Permit root and nested \pathref{AGENTS.md} files plus \pathref{.agent/} as
operational-guidance exceptions. They define durable engineering rules,
planning protocol, and reusable guidance scaffolds, but do not own current
architecture narration.\\
\textbf{Why:} One technical-documentation tree still gives implementation and
decision prose a discoverable source of truth, while scoped agent guidance must
sit where instruction inheritance applies. Treating those operational files as
ordinary architecture prose either hides local rules or duplicates the design
book. The chart pack is also a deliberate second generated view rather than a
competing prose book.\\
\textbf{Consequences:} Root guidance contains project-wide commands and
invariants; nested guidance contains only durable local specializations. Active
guidance has no template placeholders and links back to canonical documents.
Long-form architecture, current file maps, rationale, and dated history remain
under \pathref{docs/}. Generated PDFs remain outputs, not sources.\\
\textbf{Refines:} D-DOC1's single-tree and single-entrypoint wording without
rewriting its historical record.
\end{decisionbox}

%----------------------------------------------------------------------------
\subsection*{D-E6 --- Illumo library and IllumoGame product split}
\begin{decisionbox}[title={D-E6 \quad 2026-08-18}]
\textbf{Decision:} Keep one repository but establish sibling CMake projects.
\texttt{Illumo}/\texttt{Illumo::Illumo} is a static library containing the
generic host, foundation, reusable services, rendering, generic console/UI,
assets, and \symbolref{DebugModule}. \texttt{IllumoGameCore} contains cellular
simulation, rulesets, product configuration, persistence, and game commands.
\texttt{IllumoGame} owns composition, platform entry/dialogs, the
\texttt{std::chrono} frame loop, and runtime staging. Supported library headers
live under \pathref{Illumo/Include/Illumo}; OpenGL implementation headers remain
private, while \texttt{Illumo::TestSupport} exposes test-only support. Startup
is fallible, required-module failure rolls back accepted modules, and library
code does not terminate the process.\\
\textbf{Why:} The cellular simulator is now a concrete downstream consumer, so
the reusable boundary can be expressed and compiled without carrying product
policy into the library.\\
\textbf{Consequences:} Root \texttt{cmake -S .} is the canonical workspace;
\texttt{cmake -S Illumo} validates the standalone library. Independent test
runners aggregate under \texttt{IllumoWorkspace}, and combined LLVM coverage
retains the 85\% production-line gate. Illumo supplies runtime assets and
licenses through a staging helper; IllumoGame separately seeds its default
\texttt{envvars.json}. This does not add an install/export package, shared DLL,
stable binary ABI, dependency, second repository, or old Illumo executable
compatibility target.\\
\textbf{Supersedes:} D-R19's deferred-extraction rule.
\end{decisionbox}

%----------------------------------------------------------------------------
\subsection*{D-E7 --- Illumo owns the system and platform runtime}
\begin{decisionbox}[title={D-E7 \quad 2026-08-18}]
\textbf{Decision:} Illumo owns every non-game runtime concern: process entry,
native dialogs and SaveLoad contract, logger lifetime, the
\texttt{std::chrono} loop, module composition, \symbolref{SysCmdLine}, and
\symbolref{BuildInfo}. IllumoGame production code is limited to
\pathref{Source/Game} and \pathref{Source/Rulesets}. It exports an
\symbolref{IllumoApplicationDefinition} with identity, CA CLI/default metadata,
and a required-module factory. Illumo's entry invokes the generic runner; the
game supplies dialog labels, while platform code performs native interaction.\\
\textbf{Why:} Bootstrap, process lifetime, parsing mechanics, and build metadata
are engine/system responsibilities. Keeping them in the consumer would make
IllumoGame a second host rather than a game-only module.\\
\textbf{Consequences:} \texttt{IllumoPlatformEntry} is an engine-owned object
target linked into the product executable. Illumo remains independent of CA
types and ruleset validation; it sees only callbacks, option descriptors, and
\symbolref{IModule}. IllumoGame keeps product policy and the executable identity,
while supported and scaffold platform implementations live under
\pathref{Illumo/Source/Platform}.\\
\textbf{Supersedes:} D-E6 only where D-E6 assigned composition, platform
entry/dialogs, and the frame loop to IllumoGame. The sibling split, dependency
direction, test split, runtime staging, and D-N2 product identity remain in
force.
\end{decisionbox}

%----------------------------------------------------------------------------
\subsection*{D-E8 --- Persistent scene hierarchy with generational handles}
\begin{decisionbox}[title={D-E8 \quad 2026-08-19}]
\textbf{Decision:} Add a product-agnostic, main-thread-affine
\symbolref{SceneGraph} to Illumo. The graph owns node storage and exposes
graph-ID-plus-slot-plus-generation \symbolref{SceneNodeHandle} values. It
provides deterministic parent/child order, cycle rejection, subtree
destruction, local and cached world transforms, dirty propagation, and local
enabled/visible state. Each node may borrow one
\symbolref{ISceneRenderAttachment}; render traversal supplies the resolved
world transform and the attachment appends backend-neutral tokens. The graph
enters the existing per-frame \symbolref{Scene} as one drawable.\
\textbf{Why:} Illumo is intended to support future projects beyond the
cellular simulator. Rebuilding persistent object identity and hierarchy every
frame wastes CPU work and leaves ownership, transform propagation, and stale
references implicit. A small retained hierarchy supplies those contracts
without requiring an ECS or changing the renderer submission model.\
\textbf{Consequences:} The per-frame Rendering Scene remains a non-owning
layered extraction list. SceneGraph owns nodes but not attachments or backend
resources; handles from another graph or an older generation are rejected.
Traversal and subtree operations remain iterative. V1 adds no ECS components,
serialization, update callbacks, prefabs, retained UI, culling structure,
physics, scripting, or CSim cell migration. \symbolref{DebugDraw3D} is the
first render attachment and composes node world state with its existing model
matrix.\
\textbf{Refines:} D-E4. The archived raw-pointer SceneObject graph remains
superseded; only D-E4's per-frame render-list role remains current.
\end{decisionbox}

%----------------------------------------------------------------------------
\subsection*{D-N2 --- IllumoGame is the simulator identity}
\begin{decisionbox}[title={D-N2 \quad 2026-08-18}]
\textbf{Decision:} Retain \textbf{Illumo} as the repository, reusable library,
and runtime-host class name. Use \textbf{IllumoGame} for the cellular-simulator
executable, window title, CLI help/version output, console branding, product
text, game test runner, and \texttt{IllumoGame.*} test namespace. Preserve the
\texttt{.illumo} extension, save magic, and version 3/2/legacy compatibility.\\
\textbf{Why:} The visible product identity should make the new dependency
direction unambiguous without needlessly breaking persisted data.\\
\textbf{Consequences:} There is no executable named \texttt{Illumo}; library
cases remain \texttt{Illumo.*}, product cases become \texttt{IllumoGame.*}, and
format-facing compatibility terminology remains Illumo where it denotes the
existing file contract.\\
\textbf{Supersedes:} D-N1 for product, executable, and test identity.
\end{decisionbox}

%----------------------------------------------------------------------------
\subsection*{D-N4 --- CSim is the simulator product identity}
\begin{decisionbox}[title={D-N4 \quad 2026-09-16}]
\textbf{Decision:} Present the cellular simulator as \textbf{CSim} in its
window, menus, console, CLI help, native dialogs, and user-facing messages.
Use \texttt{.csim} for newly named save files. Retain \texttt{IllumoGame} for
technical executable, CMake target, source-tree, and test identifiers, and
continue loading existing \texttt{.illumo} files.\\
\textbf{Why:} CSim is the client application's product name, while Illumo is
the reusable engine/library and IllumoGame is an implementation identifier.\\
\textbf{Consequences:} Save commands and dialogs default to \texttt{.csim};
extensionless loads try \texttt{.csim} before legacy \texttt{.illumo}. Save
magic, version 4 structure, and version 3/2/dense compatibility remain
unchanged. Existing files require no migration.\\
\textbf{Supersedes:} D-N2 for product-visible simulator identity and D-N1/D-N2
for the canonical new-save extension.
\end{decisionbox}

%----------------------------------------------------------------------------
\subsection*{D-N3 --- IllEd is the world-editor identity}
\begin{decisionbox}[title={D-N3 \quad 2026-08-29}]
\textbf{Decision:} Add \textbf{IllEd} as a second in-tree Illumo application.
Use \texttt{IllEd} for the executable, window title, CLI help/version output,
console branding, editor test runner, and \texttt{IllEd.*} test namespace.
Persist authored scenes as UTF-8 JSON \texttt{.ilsc} files, not
\texttt{.illumo}.\\
\textbf{Why:} Illumo is the reusable runtime. IllumoGame is one product (the
cellular simulator). IllEd is the editor used to bootstrap later Illumo
applications, analogous to Unreal Editor versus Unreal, not a second CA
sandbox.\\
\textbf{Consequences:} Workspace CMake builds three primary executables
(\texttt{IllumoGame}, \texttt{IllEd}, plus test runners). IllEd links only
\texttt{Illumo::Illumo}. Simulator saves remain sparse version 3.
\end{decisionbox}

%----------------------------------------------------------------------------
\subsection*{D-E10 --- .ilsc is editor-owned scene interchange}
\begin{decisionbox}[title={D-E10 \quad 2026-08-29}]
\textbf{Decision:} Scene files are a product-owned document, not a
\texttt{SceneGraph} feature. \texttt{.ilsc} version 1 is pretty-printed UTF-8
JSON with \texttt{format = ilsc}, stable string node ids, parent references,
transforms, enabled/visible flags, a camera view hint, and attachment kinds
\texttt{empty} and \texttt{solid\_cube}. \texttt{SceneNodeHandle} values are
never persisted. Loads validate completely, then replace live state.\\
\textbf{Why:} D-E8 forbids serialization, names, and attachment ownership on
the graph. Handles are graph-local generations. Other tools should compose
through a readable file rather than through SceneGraph internals.\\
\textbf{Consequences:} IllEd reconstructs a \texttt{SceneGraph} from the
document and owns \texttt{MeshVisual} attachments. Extra JSON keys may be
ignored; unknown kinds and cycles fail closed. Promote the codec into Illumo
when a second loader exists. This does not change \texttt{.illumo} or CA
persistence.\\
\textbf{Refines:} D-E8.\\
\textbf{Superseded by:} D-E18 and D-E19 (2026-09-23). The codec moved into
\texttt{Illumo::Content} once every program needed it, and version 1 is no
longer read.
\end{decisionbox}

%----------------------------------------------------------------------------
\subsection*{D-E11 --- Attachment bounds and camera-relevant shadows}
\begin{decisionbox}[title={D-E11 \quad 2026-09-16}]
\textbf{Decision:} Add finite ordered \symbolref{AxisAlignedBounds3} and an
optional local-bounds query to \symbolref{ISceneRenderAttachment}.
\symbolref{SceneGraph} transforms those bounds on demand, exposes world bounds,
and rejects only attachments wholly outside the active camera frustum during
its existing iterative O(n) traversal. Unknown or invalid bounds and matrices
fail open, survivor order is stable, and a culled parent attachment does not
prune children.\
\textbf{Decision:} \symbolref{Renderer} retains directional-shadow caster
descriptors for the frame. The first camera-visible caster selects the light;
the reconstructed camera frustum is conservatively extruded toward it by the
caster's configured horizon (100 world units by default). Only intersecting
bounds contribute to the fitted shared light projection or emit depth tokens.
Invalid camera reconstruction preserves the former all-caster behavior.
\symbolref{MeshVisual} reports immediate conservative bounds for procedural,
managed-mesh, and sprite geometry; billboard sprites use a view-independent
enclosure. A frame serial resets previous-MVP history after a culling gap.\
\textbf{Why:} The graph could express hierarchy but not visibility, while the
shared shadow fit paid for and expanded around casters that could not affect
the camera-visible region. Attachment-authoritative bounds reduce token and
shadow work without changing node storage or ownership.\
\textbf{Consequences:} Direct \symbolref{MeshVisual} drawables and graph
attachments share the same Renderer-owned camera and shadow tests. Skyboxes and
attachments without reliable bounds continue rendering. No BVH, octree,
subtree-bound cache, persistence change, picking migration, or second shadow
map is introduced.\
\textbf{Refines:} D-E8, D-R21, and D-R22.
\end{decisionbox}

%----------------------------------------------------------------------------
\subsection*{D-E12 --- Compiled SceneGraph v2}
\begin{decisionbox}[title={D-E12 \quad 2026-09-17}]
\textbf{Decision:} Replace v1 node storage with authoritative SoA/TRS state,
intrusive ordered links, lazy compiled preorder, dirty watermark propagation,
revision-cached attachment and subtree bounds, and a derived binned-SAH query
index. Expose interned names, opaque payloads, ordered borrowed attachments,
handle queries, and an overflow-signaled 4,096-entry change journal. IllEd edits
this graph incrementally while retaining its file recipes and exact picking.\\
\textbf{Rationale:} Constant-time transform invalidation, stable identity,
allocation-free warmed frames, and bounded consumer updates resolve measured
v1 costs without introducing components, scene serialization, or backend policy.\\
\textbf{Consequences:} World getters remain ancestor walks; cache allocation
failure uses authoritative/linear fallbacks. Matrix setters now explicitly
lose shear through TRS decomposition. Warm indexed queries still poll bounds
revisions linearly. Generic WorkerPool remains separate from CA worker policy;
parallel scene stages remain measurement-gated. Supersedes the internal storage
and linear-only restrictions of D-E8/D-E11. See the v2 plan for validation.
\end{decisionbox}

\subsection*{D-G1 --- Pattern I/O is a side path, not save v4}
\begin{decisionbox}[title={D-G1 \quad 2026-08-21}]
\textbf{Decision:} Editor selection, in-app \texttt{CellPattern} buffers, Golly-style
RLE, Life 1.0 plaintext, OS clipboard text, and built-in stamps are a Game-owned
side path. Sparse world persistence remains version 3. RLE \texttt{rule =}
headers are ignored and never auto-switch the live ruleset.\\
\textbf{Why:} Pattern exchange is an editing convenience. Folding it into the
world save format would force a compatibility bump without changing stored
chunk topology.\\
\textbf{Consequences:} Copy writes RLE text; paste parses RLE or plaintext.
Oversize rectangles (above 256$\times$256 or 8,192 stored cells) fail the
command without mutating the grid. Finite worlds skip out-of-bounds cells
instead of wrapping stamps. Full canvas clear remains \texttt{clear\_canvas}.
\end{decisionbox}

%----------------------------------------------------------------------------
\subsection*{D-G2 --- Elementary 1D space-time neighborhood}
\begin{decisionbox}[title={D-G2 \quad 2026-08-21}]
\textbf{Decision:} \symbolref{RuleSet} distinguishes
\texttt{NeighborhoodKind::MooreCount} from \texttt{Elementary1D}. Moore rules
keep the 256$\times$9 table. Rule 90 and Rule 184 use
\texttt{nextElementary(left,center,right)} and a serial space-time
\symbolref{SparseCellGrid} advance: the source row is the maximum counted Y,
the destination is that Y plus one, and older rows remain history.\\
\textbf{Why:} Rule 184 is not a Moore popcount, and overloading the 8-neighbor
path would silently mis-evolve 1D traffic. Space-time diagrams are the
authorized presentation, not a single live row.\\
\textbf{Consequences:} Candidate, halo, frontier, and worker-pool paths stay
Moore-only. Dense \texttt{calcGeneration} remains a compatibility fixture.
\end{decisionbox}

%----------------------------------------------------------------------------
\subsection*{D-E9 --- DebugModule is a global overlay}
\begin{decisionbox}[title={D-E9 \quad 2026-08-29}]
\textbf{Decision:} Debug-build \symbolref{DebugModule} is a global overlay, not
a cell-canvas tool. The host updates optional modules before the required
product module and dispatches them afterward. \symbolref{DebugModule} consumes
Grave and open-console editing first and leaves remaining events for the
product. Product surfaces (main menu, settings, confirm dialogs, camera,
editor) yield while the console is open. Unconsumed key and character events
are discarded at the end of \symbolref{Illumo::update}. Escape closes an open
console; product modules own Escape on their own screens.\\
\textbf{Why:} The main menu drained the input queues before the overlay could
toggle the console, so Grave only worked after transitioning to the cell
canvas.\\
\textbf{Consequences:} Console, FPS, and renderer-demo remain available across
module transitions. Release still neither compiles nor registers
\symbolref{DebugModule}. The former Debug-only leftover Escape-to-close-window
path is removed; confirmed shutdown stays on the product Q/Exit surfaces.
\end{decisionbox}

%----------------------------------------------------------------------------
\subsection*{D-R21 --- Unified world look and MeshVisual}
\begin{decisionbox}[title={D-R21 \quad 2026-08-29}]
\textbf{Decision:} World objects share one look and one host:
\begin{itemize}
  \item Shape/Sprite programs (and the asset fallback) position vertices with
        \texttt{uMVP} only. Overlay chrome supplies a Y-down screen ortho;
        world draws supply camera view-projection times node world.
  \item A sprite is a textured quad. \texttt{MeshFacing::Billboard} faces the
        camera; the default is world-aligned. Orthographic vs perspective is a
        \symbolref{Camera} projection, not a second primitive system.
  \item \symbolref{MeshVisual} is the world mesh host and
        \symbolref{ISceneRenderAttachment}. Product UI remains
        \symbolref{GameVisual} painter composition (D-R20). CanvasView stays a
        world-space GameVisual sprite.
\end{itemize}
\textbf{Why:} Parallel 2D (\texttt{uUsePixels}) and 3D (private DebugDraw3D
view-projection) stacks made adding objects and composing them on
\symbolref{SceneGraph} implicit. One matrix contract is the simple, powerful
path.\\
\textbf{Consequences:} \pathref{WorldLook.h}, \pathref{MeshVisual.h},
\pathref{Camera.h}; \texttt{DebugDraw3D} is removed.
\texttt{render3dTest} drives \symbolref{MeshVisual} attachments with the product
camera. No render graph, materials, or model loader.\\
\textbf{Refines:} D-R14, D-R15. Overlay vs world remains a camera/layer split,
not a 2D engine vs 3D engine.
\end{decisionbox}

%----------------------------------------------------------------------------
\subsection*{D-R22 --- One scene-level directional shadow pass}
\begin{decisionbox}[title={D-R22 \quad 2026-09-16}]
\textbf{Decision:} Directional shadows are extracted once per rendered scene,
before color passes. Visible World drawables and
\symbolref{ISceneRenderAttachment} values contribute transformed triangle
bounds and depth tokens. \symbolref{Renderer} fits one orthographic light-space
matrix to the combined caster bounds, reuses one depth framebuffer, clears it
once, renders every caster, and exposes the same matrix and texture to all
receivers. \symbolref{SceneGraph} forwards collection, depth, and color phases
through its existing iterative enabled/visible transform traversal.
\symbolref{MeshVisual} retains geometry and per-receiver bias/filter values but
owns no shadow framebuffer.\
\textbf{Why:} A private map in each MeshVisual repeated the pass for every
object and could contain only that object's geometry, so separate objects could
not shadow one another. Fitting every private map around the origin also clipped
translated casters.\
\textbf{Consequences:} Direct World meshes and graph attachments shadow each
other in one pass. Frame-list order deterministically selects the first visible
caster's directional light; maximum requested map size, minimum radius, and
light distance cover the participating set. Existing MeshVisual tuning setters
remain source-compatible. The change adds no render graph, retained culling
structure, backend API, or persistence format.\
\textbf{Refines:} D-R14 and D-R21.
\end{decisionbox}

%----------------------------------------------------------------------------
\subsection*{D-T2 --- Workspace clang-tidy gate}
\begin{decisionbox}[title={D-T2 \quad 2026-08-30}]
\textbf{Decision:} First-party C++ is linted by workspace \texttt{clang-tidy}.
\texttt{ILLUMO\_ENABLE\_CLANG\_TIDY=ON} adds \texttt{IllumoTidy}, which reads a
Ninja/Clang \pathref{compile\_commands.json} and runs LLVM \texttt{clang-tidy}
with \pathref{.clang-tidy}. Diagnostics are errors.
\texttt{python build.py tidy} is the supported front end. The default MSVC
workspace compile does not invoke clang-tidy. Combined coverage remains a
separate Clang/LLVM tree (D-T1).\\
\textbf{Why:} Manual one-off clang-tidy invocations were not a repeatable gate,
and wrapping every MSVC compile would double build time and mix Clang's parser
with MSVC flags.\\
\textbf{Consequences:} Vendored third-party translation units are excluded.
\texttt{modernize-use-auto} stays disabled because house style forbids
\texttt{auto} (D-008). Standalone \texttt{cmake -S Illumo} rejects the tidy
option the same way it rejects combined coverage.
\end{decisionbox}

%----------------------------------------------------------------------------
\subsection*{D-T3 --- clang-tidy runs during first-party compiles}
\begin{decisionbox}[title={D-T3 \quad 2026-08-30}]
\textbf{Decision:} \texttt{ILLUMO\_ENABLE\_CLANG\_TIDY} defaults to ON.
First-party targets receive \texttt{CXX\_CLANG\_TIDY} and are linted during
the ordinary compile with \pathref{.clang-tidy}. Diagnostics remain errors.
Disable with \texttt{-DILLUMO\_ENABLE\_CLANG\_TIDY=OFF} or
\texttt{python build.py build --no-tidy}. The batch \texttt{IllumoTidy} target
and \texttt{python build.py tidy} remain available. Combined coverage still
disables tidy so the instrumented tree stays a coverage-only configure
(D-T1).\\
\textbf{Why:} An opt-in-only Ninja target was not part of the default build
the workspace actually runs.\\
\textbf{Consequences:} LLVM \texttt{clang-tidy} must be on \texttt{PATH} unless
the option is turned off. Vendored sources (including TracyClient) are not
linted. Standalone \texttt{cmake -S Illumo} lints first-party library sources
the same way; it still rejects combined coverage.\\
\textbf{Supersedes:} D-T2's default-MSVC-does-not-run and standalone-reject
details.
\end{decisionbox}

\subsection*{D-R22 --- Cubemap and SkyboxVisual rendering}
\begin{decisionbox}[title={D-R22 \quad 2026-09-04}]
\textbf{Decision:} The rendering subsystem supports 3D environment skyboxes via
cubemaps:
\begin{itemize}
  \item \symbolref{IBackend} gains \texttt{CreateCubemap} allocating a
        generational \symbolref{TextureHandle} bound to
        \texttt{GL\_TEXTURE\_CUBE\_MAP} with 6 face uploads, seamless filtering,
        and \texttt{GL\_CLAMP\_TO\_EDGE} wrap.
  \item \symbolref{AssetManager} loads individual face files via
        \texttt{acquireCubemap} and automatically extracts 6 faces from horizontal
        crosses (4:3), vertical crosses (3:4), and strips (6:1) via
        \texttt{acquireCubemapFromCross}.
  \item Built-in style \texttt{RenderStyleId::Skybox} renders a unit cube
        at the far depth plane using \texttt{pos.xyww} projection,
        \texttt{GL\_LEQUAL} depth testing, and no backface culling.
  \item \symbolref{SkyboxVisual} is the world cubemap host and
        \symbolref{ISceneRenderAttachment}, emitting tokens with a
        translation-stripped camera view-projection matrix. Consuming
        viewers (e.g. \texttt{IllMeshViewer}) expose live toggles and runtime
        configuration.
\end{itemize}
\textbf{Why:} 3D scene visualization and model viewing require environment
context without polluting the backend-neutral token queue or breaking the
single-pipeline render model.\\
\textbf{Consequences:} \pathref{SkyboxVisual.h}, \pathref{GLTexture.h},
\pathref{IBackend.h}, \pathref{AssetManager.h}; tested via headless mock tokens
and E2E suite; integrated into \texttt{IllMeshViewer}.
\end{decisionbox}

%----------------------------------------------------------------------------
\subsection*{D-R23 --- Pixel clipping and conservative quad culling}
\begin{decisionbox}[title={D-R23 \quad 2026-09-15}]
\textbf{Decision:} During pixel-space rebuilds, \symbolref{GameVisual} tests
each final quad against the logical viewport and omits wholly excluded geometry
before buffer upload. An optional top-left logical-pixel clip further narrows
that rectangle. Partial overlap emits ordinary geometry bracketed by
\texttt{Renderer::pushClipRect}/\texttt{popClipRect}; the renderer intersects
nested physical-pixel scissors and restores the prior state.\\
\textbf{Why:} Offscreen UI, sprites, and glyphs previously consumed dynamic
quad capacity, CPU vertex writes, upload bytes, and vertex work even though the
viewport or a product clip discarded every fragment. Existing low-level
scissor tokens could clip pixels but did not remove that work or safely nest.\\
\textbf{Consequences:} Painter order, adjacent-only batching, partial-primitive
rasterization, and explicit low-level scissor emission are preserved. Window
or UI-scale changes invalidate the cached pixel cull rectangle and rebuild
visibility. World-space \symbolref{GameVisual}, \symbolref{MeshVisual}, and
SceneGraph extraction are unchanged; no frustum hierarchy or spatial index is
introduced.
\end{decisionbox}

%----------------------------------------------------------------------------
\subsection*{D-R24 --- Shared immutable mesh assets}
\begin{decisionbox}[title={D-R24 \quad 2026-09-16}]
\textbf{Decision:} Extend \symbolref{AssetManager} with reference-counted
immutable static meshes. Canonical file paths plus geometry-affecting
\texttt{MeshLoadOptions} identify cached uploads; in-memory \texttt{MeshData}
can be enrolled once and explicitly retained. \symbolref{MeshVisual} borrows the
managed \texttt{MeshHandle} and immutable draw count while keeping only
per-instance transform, tint, lighting, shadow, and motion state. Procedural
lines, shapes, and sprites retain their visual-owned dynamic buffers.\\
\textbf{Why:} Copying imported vertices into every \symbolref{MeshVisual}
duplicated CPU storage and GPU uploads for nodes that displayed the same model.
Backend generational mesh handles already supplied the stable binding needed to
separate resource lifetime from instance state.\\
\textbf{Consequences:} Two nodes can draw one uploaded mesh with independent
world transforms and tint. Final manager release destroys the backend resource;
visual destruction does not. Mesh loading and enrollment are synchronous and
main-thread affine in this milestone. Materials, mesh hot reload, and automatic
instance aggregation remain deferred until a measured workload justifies their
additional contracts.\\
\textbf{Refines:} D-R17's earlier texture/shader-only managed scope and D-R21's
world-mesh host.
\end{decisionbox}

%----------------------------------------------------------------------------
\subsection*{D-R25 --- Immutable scene snapshot boundary}
\begin{decisionbox}[title={D-R25 \quad 2026-09-17}]
\textbf{Decision:} SceneGraphDrawable consumes one graph-owned immutable view
per renderer frame for collection, depth, and color. Two retained buffers carry
ordered matrices, bounds, node handles, and borrowed attachment pointers.
Validate view lifetime before every callback; stale views report frame errors.
Renderer/token/backend contracts and main-thread execution remain unchanged.\\
\textbf{Rationale:} Rendering no longer traverses mutable hierarchy three times.
Explicit invalidation separates graph mutation from borrowed-content lifetime.\\
\textbf{Consequences:} Ordinary edits affect the next snapshot; attachment
changes, destruction, ring reuse, and teardown retire views. Already emitted
payloads still outlive submission. Shared light fitting follows collection,
so shadow relevance uses snapshot bounds in the depth pass; off-camera casters
remain present. A one-release sceneSnapshotExtraction switch retains direct
traversal. Recorded payloads and GL-context migration require a separate decision.
\end{decisionbox}

\subsection*{D-CAP1 --- Bounded standalone rendering and charter baseline}
\begin{decisionbox}[title={D-CAP1 \quad 2026-09-11}]
\textbf{Decision:} Reuse the token renderer for direct and SceneGraph frame capture
in a hidden OpenGL context. Capture owns a bounded service lifetime, not an
application loop. Readback and strict diagnostics are public contracts; PNG
publication does not replace existing artifacts.
\textbf{Why:} The owner authorized the first two charter milestones: reliable
baseline behavior and independently usable rendering tools.
\textbf{Consequences:} Input IDs are validated, retired and never reused within
a manager; generated copies record commit/dirty/unknown provenance. Broader 3D
framework direction is separate from current code. Physics, object destruction,
shared ticks, additional backends and Linux support remain later work. Existing
application and persistence contracts remain. See the milestone execution plan.
\end{decisionbox}

\subsection*{D-L1 --- Product close negotiation}
\begin{decisionbox}[title={D-L1 \quad 2026-09-13}]
\textbf{Decision:} The runner explicitly negotiates a pending native close with
started modules through optional \texttt{OnCloseRequested}. Default acceptance
preserves existing consumers; deferral clears the window flag and keeps frames
running. Required startup/transition failure remains terminal.
\textbf{Why:} Native close previously bypassed IllEd's unsaved-document guard.
\textbf{Consequences:} IllEd routes native and toolbar exit through its existing
Save/Discard/Cancel confirmation. Approved discard may close a dirty document.
Custom windows used with deferring modules implement close cancellation. No
editor policy moves into Illumo.
\end{decisionbox}

\subsection*{D-IO1 --- Finalized save publication}
\begin{decisionbox}[title={D-IO1 \quad 2026-09-13}]
\textbf{Decision:} Product saves use the shared platform atomic-file operation:
exclusive sibling staging, writer/flush/close checks, then replacement.
\textbf{Why:} Direct truncation destroyed prior valid saves before late I/O
errors could be reported.
\textbf{Consequences:} Product formats and encoders remain product-owned.
Failure leaves prior destination bytes intact. This is not a power-loss
durability or security-metadata preservation guarantee. Windows is verified;
POSIX rename remains unverified scaffolding.
\end{decisionbox}

\subsection*{D-RP1 --- Host defaults and application pass precedence}
\begin{decisionbox}[title={D-RP1 \quad 2026-09-13}]
\textbf{Decision:} Scene retains host default pass lists separately from
application overrides. Nonempty overrides take precedence; clearing them
restores the current fallback. Host motion blur writes only defaults.
\textbf{Why:} Per-frame host configuration previously erased passes installed
in module Start or Update.
\textbf{Consequences:} Existing application setters remain available. Empty
lists and reset restore defaults. Effective pass queries include fallbacks.
Resource ownership and synchronous token submission remain unchanged.
\end{decisionbox}

\subsection*{D-GC1 --- Product-owned rules catalog discovery}
\begin{decisionbox}[title={D-GC1 \quad 2026-09-14}]
\textbf{Decision:} Game's RuleCatalogLoader owns catalog file reads and ordered
discovery. The required-module factory loads the singleton once before either
menu or direct-game construction. RuleSetRegistry constructs built-ins and
accepts catalog text with transactional validation.\\
\textbf{Why:} Domain registry construction previously called Win32 and read
ambient files, contradicting the Rulesets dependency boundary.\\
\textbf{Consequences:} Executable/current/product-directory precedence and
first-valid fallback remain. Product C++ file-loading calls migrate to the
loader; independent registries require explicit catalog loading. No engine
API or catalog format changes. Native path discovery stays engine-owned.
\end{decisionbox}

\subsection*{D-GC2 --- Data-driven rules and runtime workshop}
\begin{decisionbox}[title={D-GC2 \quad 2026-09-15}]
\textbf{Decision:} Shipped and custom rules use versioned JSON definitions
compiled into data-backed rules. Keep the schema-v1 reader and stable rule IDs;
layer the working-directory \pathref{rulesets.user.json} after the selected
base catalog. F2 edits a staged custom rule and Save \& Apply drains the
simulation, validates live states, atomically persists the draft, then
activates it.\\
\textbf{Supersedes:} D-GC1's built-in-construction clause; its file-discovery
and Rulesets dependency-boundary clauses remain in force.\\
\textbf{Why:} Ruleset behavior and palettes should have one data source, and
specific rule experiments should not require recompiling or silently changing
IDs already stored in worlds.\\
\textbf{Consequences:} Life-like, Generations, Moore-table, and elementary 1D
families compile through the existing simulator contracts. The sparse-domain
background/counting invariants and \texttt{.illumo} v3 format stay unchanged.
Moore tables remain JSON-edited; the in-game menu focuses on common family
parameters and palettes. A custom world still requires its rule catalog at
load time.
\end{decisionbox}

\subsection*{D-GC3 --- Separate ruleset identity from family}
\begin{decisionbox}[title={D-GC3 \quad 2026-09-15}]
\textbf{Decision:} Keep each ruleset's stable ID and display name independent
from a typed \texttt{RuleFamily}. The Ruleset Workshop exposes starter rule,
family, and ruleset name as separate choices. Changing family retains the
custom ID/name and loads that family's starter parameters.\
\textbf{Why:} A family describes transition semantics; it should not be
conflated with the named, save-referenced rule that uses those semantics.\
\textbf{Consequences:} Schema-v1/v2 JSON family strings and existing
\texttt{.illumo} rule IDs remain unchanged. Family conversion is explicit at
the registry boundary and the workshop rebuilds controls for the selected
family.\
\end{decisionbox}

\subsection*{D-GC4 --- Family-owned cell schemas and pair persistence}
\begin{decisionbox}[title={D-GC4 \quad 2026-09-15}]
\textbf{Decision:} Store family and ruleset as separate data definitions.
Families own the cell-state schema, labels, colors, and simulation model;
every ruleset references exactly one family and owns its transition behavior.
F1, New Simulation, runtime status, and saves carry the validated pair.
Write \texttt{.illumo} version 4 with both IDs and retain readers that derive
the family from the rule in older saves. Layer
\pathref{families.user.json} and \pathref{rulesets.user.json} together after
the selected valid base catalog pair.\\
\textbf{Supersedes:} D-GC2's single rules overlay and version-3 save outcome;
D-GC3's separation of UI identity is retained and now has separate family data.\\
\textbf{Why:} Cell appearance and available states are family properties;
transition parameters belong to one named ruleset bound to that family. This
keeps cell schemas reusable while preventing meaningless cross-family pairs.\\
\textbf{Consequences:} Legacy catalogs normalize in memory, old save formats
remain readable, and pair validation occurs before catalog or world mutation.
Built-in family edits become custom copies; user families are persisted before
the rules that reference them.
\end{decisionbox}

\subsection*{D-GC5 --- Isolated cyclic multistate interaction}
\begin{decisionbox}[title={D-GC5 \quad 2026-09-15}]
\textbf{Decision:} Add a typed \texttt{cyclic} family whose rules advance a
cell by a nonzero coprime step when a threshold number of Moore neighbors hold
that successor state. Expose full 256-state neighbor counts through
\texttt{RuleSet}; evaluate them in a separate serial sparse kernel that reuses
transactional output and change journals.\\
\textbf{Why:} Count-of-state-zero and Moore-table models cannot express many
cell kinds directly chasing one another. A general pairwise language would be
harder to author and unnecessary for the requested dynamic ecology.\\
\textbf{Consequences:} Existing optimized binary/counting paths and save
formats remain unchanged. State 1 stays sparse background and is quiescent
without successor neighbors. Cyclic rule JSON owns threshold and step; the
workshop edits both. The first histogram kernel prioritizes determinism and
correctness over parity with specialized sparse performance.
\end{decisionbox}

\subsection*{D-GC6 --- Researched interaction and extended-range families}
\begin{decisionbox}[title={D-GC6 \quad 2026-09-15}]
\textbf{Decision:} Add typed \texttt{species\_life} and
\texttt{larger\_than\_life} families. Colorized Life uses the existing
full-state Moore histogram to preserve or select a live species.
Larger-than-Life adds a bounded extended-range contract with square/circular
shape, optional center counting, and inclusive transition intervals; its first
sparse implementation is an isolated serial correctness kernel. Store an
optional deterministic starter strategy with each rule.\\
\textbf{Why:} Documented multi-species, long-decay, cyclic, and wide-range
systems provide visibly different interactions, but a Life glider is not a
meaningful initial condition for most of them. Repeated radius-one updates are
not equivalent to a wide-neighborhood rule.\\
\textbf{Consequences:} Immigration, QuadLife, Bosco/Bugs, Bugsmovie, Globe,
Banners, Transers, Fireworks, and classic 14-color CCA ship as data. Existing
catalog files remain readable because the new keys are additive. B0 remains
invalid for sparse infinite worlds; extended range is capped at 16 and does not
use the specialized radius-one candidate, memo, or worker paths. Starter soups
are coordinate-hashed and repeatable. Continuous and block-partitioned
automata remain out of scope.
\end{decisionbox}

\subsection*{D-GC7 --- Directional agents, lattice gas, and multilevel chemistry}
\begin{decisionbox}[title={D-GC7 \quad 2026-09-15}]
\textbf{Decision:} Add typed \texttt{hodgepodge}, \texttt{turmite},
\texttt{lattice\_gas}, and \texttt{dominance} families. Preserve cardinal
direction through a four-entry von Neumann rule contract and an isolated serial
sparse evaluator. Reuse the complete Moore histogram for chemical levels and
species predators. Keep state 1 as sparse background through explicit state
encoding maps.\\
\textbf{Why:} Unordered live counts cannot express moving directional agents,
particle streaming, or weighted chemical infection. These models add genuinely
different interactions while still fitting the byte-state synchronous domain.\\
\textbf{Consequences:} RL/RRL/RRLL Turmites, 16-state HPP gas, 101-level
Hodgepodge, and five-species RPSLS ship as data. Directional and histogram
kernels publish transactionally but initially remain serial. Multi-agent
Turmite collisions annihilate deterministically. The synchronous RPSLS rule is
inspired by, not identical to, stochastic Monte Carlo scheduling. FHP/hexagonal
gas, continuous chemistry, and generic Turing-machine controllers remain out of
scope.
\end{decisionbox}

\subsection*{D-GC8 --- Rule tables, sandpiles, long-range cycles, and trails}
\begin{decisionbox}[title={D-GC8 \quad 2026-09-23}]
\textbf{Decision:} Add typed \texttt{von\_neumann\_table} and
\texttt{sandpile} families on the existing directional contract. Rule tables
are stored as Golly \texttt{@TABLE} text (variables, comma or compact
transitions, \texttt{none}, \texttt{reflect\_horizontal},
\texttt{rotate4}, \texttt{rotate4reflect}, or \texttt{permute} symmetry) and
compile to a dense $n^5$ lookup for at most twelve states with Golly
semantics: first match wins, repeated variables bind, unmatched neighborhoods
keep their center, and Golly state 0 swaps with Illumo background state 1.
Sandpiles use heights 0--7 plus a pulsed grain source whose phase count comes
from the family. Cyclic rules gain an optional range and square, circular, or
diamond shape and an optional inert background that is not a phase of the
cycle; Larger than Life gains the diamond shape and Golly's C-state
decay trail. Extended-range kernels count the state named by
\texttt{RuleSet::getExtendedCountedState} and read a per-target halo window.
Rules may carry an exact Golly-RLE starter.\\
\textbf{Why:} Self-replicating loops, sandpile fractals, Griffeath spirals, and
trailing Larger-than-Life spaceships are well-documented behaviors that the
unordered-count and radius-one cyclic contracts could not express. Storing
published tables verbatim keeps them auditable against their sources.\\
\textbf{Consequences:} Langton, Byl, Chou--Reggia, SDSR, and Evoloop loops,
three sandpile starters, six Griffeath rules, four trailing LtL rules, and 45
classic Life-like and Generations rules ship as data. Loop evolution is checked
against an independent reference implementation. Radius-one square cyclic rules
keep the Moore histogram path and serialized form. A full cycle advances the
empty background too, so long-range soups would grow without bound on the
infinite canvas; the shipped Griffeath rules use the inert background, while
existing full-cycle rules are unchanged. The halo window replaces
per-neighbor hash lookups for extended and directional rules without changing
their results; the kernels remain serial and use the existing transactional
publication. Moore-neighborhood tables, JvN29-scale state counts, and
Margolus partitioning remain out of scope.\\
\textbf{Extends:} D-GC6, D-GC7
\end{decisionbox}

\subsection*{D-GC9 --- Lenia as quantized weighted-kernel rules}
\begin{decisionbox}[title={D-GC9 \quad 2026-09-25}]
\textbf{Decision:} Add a typed \texttt{lenia} family model on a new
\texttt{WeightedKernel} neighborhood contract instead of a floating-point cell
domain. Cells stay bytes: a family of $n \le 256$ states holds intensity
levels $0..n-1$ (background state 1 is level 0, states $2..n-1$ are levels
$1..n-2$, state 0 is the full level). Rules carry Lenia's radius $R \le 16$,
time steps $T$, growth centre $\mu$ and width $\sigma$, one to four kernel
ring peaks $\beta$, and polynomial, exponential, or step core and growth
shapes. Compilation quantizes the normalized kernel to integer tap weights and
$G(u)/T$ to a fixed-point level delta per potential bin, evaluating
exponentials with a deterministic series rather than the platform
\texttt{exp}. The evaluator sums weight $\times$ level exactly in integers and
rounds the clipped update to the nearest level, which is Lenia's own
$P$-quantized update with $P = n-1$.\\
\textbf{Why:} Lenia's continuous state fits the sparse byte domain once it is
quantized, and Chan defines state resolution $P$ as part of the model. Keeping
bytes reuses storage, saves, presentation, editing, and simulation lanes
unchanged. Exact integer potentials and platform-independent tables keep the
control store, WASM lanes, and native oracle bit-identical, which lane parity
requires.\\
\textbf{Consequences:} A 256-level family and six species (Orbium, Orbium
bicaudatus, Gyrorbium, Scutium, Discutium, Paraptera) ship with
\texttt{lenia\_rle} starters from Chan's MIT-licensed catalogue; Orbium's
quantized evolution matches an independent numpy model exactly and glides.
Growth that raises empty space is rejected so the infinite canvas stays sparse.
The kernel is serial and direct ($O(R^2)$ taps per cell); larger worlds rely
on simulation lanes. Species with $R > 16$, FFT convolution, multi-channel and
multi-kernel Lenia, and exact floating-point states remain out of scope.\\
\textbf{Extends:} D-GC6, D-GC8
\end{decisionbox}

\subsection*{D-GC10 --- Every cell family on the worker pool}
\begin{decisionbox}[title={D-GC10 \quad 2026-09-25}]
\textbf{Decision:} The full-state histogram, directional, extended-range, and
weighted-kernel models share one chunk-parallel driver on the grid's worker
pool. Target discovery is serial; each target chunk is evaluated from its own
per-worker halo window (the grid-wide window is gone); the change journal and
inactive-map publication run serially in target order. Work above 131,072
neighbor reads (target cells times per-cell kernel work) uses up to eight
workers, so wide kernels parallelize from two targets. Elementary 1D computes
its next row in bounded blocks split across the pool and writes it serially.
The pool gains a generic indexed-job entry point with stable per-thread slots.\\
\textbf{Why:} Only the radius-one count-table paths used the pool, leaving
colorized Life, cyclic, Hodgepodge, dominance, Turmites, HPP, rule tables,
sandpiles, Larger than Life, and Lenia on one thread. Their targets were already
independent reads of a frozen source, so parallel evaluation needs no change to
rule contracts or results.\\
\textbf{Consequences:} Every shipped rule is checked to publish byte-identical
generations with one and four workers on infinite and toroidal worlds, and each
neighborhood kind is checked to dispatch more than one worker. The histogram
kernel now reads a one-cell halo window instead of hashing each neighbor. The
WASM guest still compiles the pool serially; in the package, simulation lanes
remain the parallelism (for every rule since D-E29). The optimized
radius-one candidate, frontier, and memo paths are unchanged.\\
\textbf{Extends:} D-GC5, D-GC6, D-GC8, D-GC9
\end{decisionbox}

\subsection*{D-DEP1 --- Narrow public math dependency exposure}
\begin{decisionbox}[title={D-DEP1 \quad 2026-09-14}]
\textbf{Decision:} Illumo exports its public headers and a CMake-managed GLM-only
include directory. Other vendor roots remain private.\\
\textbf{Why:} Public math types intentionally expose GLM, but exporting the
entire thirdparty parent directory exposed unrelated libraries too.\\
\textbf{Consequences:} GLM header paths and types remain unchanged. Configure
tracks header content and membership; deleted headers leave the generated
include tree. Vendored sources and versions are unchanged. Consumers of other
vendor APIs must declare those dependencies explicitly.
\end{decisionbox}

\subsection*{D-PROF1 --- Bounded in-game frame timing pie}
\begin{decisionbox}[title={D-PROF1 \quad 2026-09-14}]
\textbf{Decision:} Illumo owns a fixed 120-frame main-thread elapsed-time
collector. DebugModule borrows it explicitly and renders a navigable GameVisual
pie in Debug and RelWithDebInfo.\\
\textbf{Why:} In-game attribution exposes expensive frame phases without
requiring a separate profiler client. CPU command submission and presentation
must be separate so pacing is not mistaken for rendering work.\\
\textbf{Consequences:} No GPU timing or asynchronous worker totals are included.
The collector is off by default; enabling does not depend on FPS, memory, or
Tracy. Number-key capture covers both events and polling without synthetic
releases. Existing context ownership, module order, and rendering tokens remain.
\end{decisionbox}

\subsection*{D-UI1 --- One shared menu shell for primitive-composed overlays}
\begin{decisionbox}[title={D-UI1 \quad 2026-09-16}]
\textbf{Decision:} Overlay behavior that more than one screen needs lives in
\symbolref{Illumo/Gui/GuiMenuShell}: \symbolref{GuiEasing} curves,
\symbolref{GuiMenuAnimator} reveal/row-stagger/selection/value-pulse/ambient/
caret clocks with reduced motion, \symbolref{GuiPanelLayout} virtual-resolution
fitting plus row-window and wheel arithmetic, and \symbolref{GuiPointerTracker}
virtual-space pointer sampling with press-edge detection. IllumoGame's menus
compose it instead of carrying private copies.\\
\textbf{Why:} Drawing and color were already shared through GuiKit and UiTheme,
but each menu had re-derived the same clocks, easing curve, UI-scale fit,
scroll clamping, and mouse-edge bookkeeping. Adding another interface meant
copying that machinery again, and the copies had already drifted in their
scale floor and reveal durations.\\
\textbf{Consequences:} A new screen supplies only its rows, layout constants,
and GameVisual composition. Menu timings are defined once, so theming and
animation stay consistent across screens. NewSimulationMenu adopts the shared
0.25 scale floor and the shared panel reveal in place of its own 0.35\,s
entrance; its per-row focus weights remain screen-local. MainMenuModule keeps
its slower title-screen entrance clock. GuiDialog keeps its own dialog-fitted
layout and clocks and shares only the easing curve. No retained widget tree is
introduced, and GuiMenuShell performs no drawing. IllEd's docked panels and
IllMeshViewer's overlay adopt the same shell for pointer edges, easing, and
viewport resolution: \symbolref{GuiPanelLayout::viewport} anchors docked chrome
to the real window where \symbolref{GuiPanelLayout::fit} would wrongly impose
the centered-overlay design floor. EditorModule keeps its own world-picking and
gizmo drag state, which is world input rather than panel chrome.
\end{decisionbox}

\subsection*{D-R26 --- Gradient quads and living-glass menu chrome}
\begin{decisionbox}[title={D-R26 \quad 2026-09-22}]
\textbf{Decision:} \symbolref{GameVisual} gains
\texttt{ShapeKind::GradientQuad}: convex, fan-ordered quads with one color per
vertex, carried by the existing shape vertex format and shader. GuiKit composes
soft chrome from them (glass panels, gradient surfaces, rounded bands for
outlines, soft shadows and glows, radial glows, vignettes, sheens, keycaps).
\symbolref{GuiMenuShell} adds \symbolref{GuiSpring} and
\symbolref{GuiSpringArray} (closed-form damped springs), \texttt{outBack} and
\texttt{inOutCubic} curves, and animator clocks for a stretching selection
span, arrival sheen, press pulse and directional value pulse. Every IllumoGame
menu and in-game overlay adopts the resulting ``living glass'' language, and
the title screen runs a live Immigration Life world behind its panel.\\
\textbf{Why:} The menus were tasteful but static: flat cards, a flat selection
fill, and a decorative background hidden under an opaque fill. Soft light and
physical motion need per-vertex color and springs; both fit the existing
primitive stream without a new shader, blur pass, or widget tree.\\
\textbf{Consequences:} Rounded rectangles pack two corner wedges per quad
(27 quads become 15, same pixels, same three-rectangle prefix). Fades target
the same color at zero alpha. Flat GuiKit helpers and GuiDialog's flat style
are unchanged. Feedback animation never delays an action and never moves a hit
area; overlays still close immediately and emit no commands when closed.
Reduced motion snaps every spring and disables decorative motion. The title
screen restores the saved ruleset preferences its \symbolref{CellContext}
would overwrite and refreshes canvas targets after each advance. Runtime
captures gain \texttt{--capture-script} and are saved opaque. No persistence,
ABI, or frame-schema change.
\end{decisionbox}

\subsection*{Template for new entries}
\begin{lstlisting}[language={}]
\subsection*{D-0XX --- short title}
\begin{decisionbox}[title={D-0XX \quad YYYY-MM-DD}]
\textbf{Decision:} ...
\textbf{Why:} ...
\textbf{Consequences:} ...
\textbf{Supersedes:} D-0YY (optional)
\end{decisionbox}
\end{lstlisting}

\subsection*{D-AUD1 --- Report contract corrections (2026-09-19)}
D-R14 is refined by D-R22/D-RP1: ordered layers coexist with the shared shadow
pass and configured offscreen/post passes. D-B1 includes Debug and RelWithDebInfo;
Release excludes DebugModule. Native GPU identifiers are private GL implementation
state; supported consumers use generational handles. D-C3's storage-independent
rules boundary is completed by moving dense evaluation and reference adapters
into test-only sources. No persistence format changes accompany these corrections.
Renderer-bound primitives reject foreign lifetimes; rejected tokens or truncated
geometry prevent frame presentation. See the report remediation plan for tests.

\subsection*{D-E13 --- IllumoGame ships only as a WASM package}
\begin{decisionbox}[title={D-E13 \quad 2026-09-22}]
\textbf{Decision:} The complete IllumoGame product executes inside
\texttt{IllumoGame.wasm}, hosted by the generic \texttt{IllumoRuntime.exe}
(Windows x64, pinned Wasmtime 48.0.2). No native game executable is built.
Inside the store, \symbolref{GuestModuleApplication} composes a guest-local
\symbolref{IllumoContext} (window snapshot, input, storage-backed settings,
camera, recording renderer, scene, console bridge, module host) so the
existing modules, menus and workshop run unchanged. Product I/O goes through
the product-owned \symbolref{CSimPlatform} seam; the simulation runner is
serial in the guest. Frame schema version 2 adds the world camera, line and
depth-tested batches, lit meshes and shadow casters so the 3D diagnostic
renders through the host's shared shadow pass.\\
\textbf{Why:} The design goal is a generic host that plays isolated game
packages; keeping a native IllumoGame alongside the package made the host a
second copy of the product.\\
\textbf{Consequences:} \texttt{IllumoGameCore} remains only as the native test
oracle for the shared sources. The default Windows build requires the pinned
WASM toolchain; other hosts build no playable game until a Linux runtime pin
exists. Host console built-ins address runtime settings, while game commands
are trampolined into the package. Generations run inside the control call
until a worker-store offload with asynchronous drains lands. The earlier
rejection of compiling \texttt{CellGameModule} into the reactor is superseded
because every context member is now a guest object.\\
\textbf{Supersedes:} D-E7 (IllumoGame application seam)
\end{decisionbox}

\subsection*{D-E14 --- Every client program is a WASM package}
\begin{decisionbox}[title={D-E14 \quad 2026-09-22}]
\textbf{Decision:} Every interactive Illumo client program runs as a WASM
package inside the one generic \texttt{IllumoRuntime.exe}. Packages live in
\texttt{apps/<name>/} beside the runtime with an \texttt{app.json} manifest
(\texttt{id}, \texttt{module}, optional \texttt{worker}, \texttt{title},
\texttt{launchAccess} \texttt{read} or \texttt{edit}, requested budgets
clamped to host ceilings) and private storage in \texttt{storage/<id>/}. The
installed apps are \texttt{game} (\texttt{IllumoGame.wasm}), \texttt{illed}
(\texttt{IllEd.wasm}) and \texttt{meshviewer} (\texttt{IllMeshViewer.wasm});
the former \texttt{game/} directory and \texttt{game.json} are replaced.
\texttt{--app} selects a package, \texttt{--open} grants one launch document
whose base name alone reaches the guest (writable when the manifest says
\texttt{edit}), and \texttt{--capture} renders to a chosen frame, writes a new
PNG from the presented backbuffer and prints one JSON result, folding the
standalone \texttt{IllumoCapture} executable into the runtime. Frame schema
version 3 adds retained host meshes (\texttt{CreateMesh},
\texttt{WriteMesh}, \texttt{ReleaseMesh}), cubemaps with \texttt{Skybox}
batches, and a per-batch blend flag. IllEd and IllMeshViewer reach files and
dialogs through the product-owned \texttt{IllEdPlatform} and
\texttt{MeshViewerPlatform} seams.\\
\textbf{Refined by:} D-E22 (2026-09-23): \texttt{illumo.json} replaces
\texttt{app.json}, and packages may be \texttt{.ilpk} archives.\\
\textbf{Why:} D-E13 proved the generic host with the game; keeping native
editor, viewer and capture executables beside it left three programs outside
the sandbox and a second capture path that did not render real applications.\\
\textbf{Consequences:} \texttt{IllEd.exe}, \texttt{IllMeshViewer.exe} and
\texttt{IllumoCapture.exe} are no longer built; \texttt{IllEdCore} and
\texttt{IllMeshViewerCore} remain as native test oracles with
\texttt{IllEdTests} and \texttt{IllMeshViewerTests}. The engine
\symbolref{FrameCapture} API and the explicit \texttt{IllumoCaptureGpuTests}
target remain; \texttt{tools/verify\_capture.py} drives
\texttt{IllumoRuntime --capture}. IllEd save, open and close confirmation
become asynchronous over opaque locations; \texttt{.ilsc} is unchanged.
Motion blur and MSAA stay host policy. Linux, lacking a runtime pin, has no
runnable applications and builds only the engine libraries and native test
suites.\\
\textbf{Extends:} D-E13 (IllumoGame WASM package)
\end{decisionbox}

\subsection*{D-E15 --- Host-chosen WASM engine modes}
\begin{decisionbox}[title={D-E15 \quad 2026-09-22}]
\textbf{Decision:} Code generation depends on \texttt{WasmEngineOptions}
(fuel metering, explicit bounds checks). The host chooses them per store and
passes them to the isolated compiler, so an artifact always matches the engine
that deserializes it; a mismatch fails closed. Explicit guest bounds checks
follow AddressSanitizer (\texttt{ILLUMO\_ENABLE\_ASAN}, default on for Debug),
not \texttt{\_DEBUG}. Package manifests choose \texttt{metering}
\texttt{fuel} (default) or \texttt{epoch}; the three first-party packages are
epoch-only, mods stay fuel-metered, and \texttt{--fuel} forces metering for a
launch. \texttt{build.py} adds the \texttt{dev} (RelWithDebInfo) and
\texttt{debug-noasan} profiles beside the \texttt{debug} sanitizer profile.\\
\textbf{Why:} Players used the sanitizer profile, where every guest memory
access carried an explicit check, and fuel instrumented every compiled block
although the epoch deadline already bounds each call.\\
\textbf{Consequences:} Epoch-only stores are bounded by their deadline alone;
fuel remains an opt-in and a comparison tool. Measured: removing fuel gains
3--7\% on the game's own generations; the Debug sanitizer profile costs
1.3--1.6$\times$ in headless play.\\
\textbf{Extends:} D-E13, D-E14
\end{decisionbox}

\subsection*{D-E16 --- Frame schema v4: retained dynamic meshes}
\begin{decisionbox}[title={D-E16 \quad 2026-09-22}]
\textbf{Decision:} Dynamic 2D meshes (Shape, Sprite, Canvas; at most 4\,MiB
per buffer) are retained on the host. \texttt{CreateMesh} with the dynamic
flag returns a ready, zero-filled mesh; a frame's trailing
\texttt{meshWrites} section patches it before any of the frame's batches draw,
and the host validates each write into a CPU shadow and uploads only dirty
spans. The guest sends only the span that differs from its mirror, draws
inline until the host copy exists, and makes a mesh inline for good (converting
that frame's by-reference draws) when it changes after being drawn in the same
frame. Remaining inline batches use one grow-only host mesh pool per style.
Versions 1--3 stay valid.\\
\textbf{Why:} Every UI, text and GameVisual buffer was re-extracted,
serialized, decoded and re-uploaded each frame, and batch reordering rebuilt
host meshes.\\
\textbf{Consequences:} The game menu's per-frame module work fell from
2.75\,ms to 0.64\,ms, IllEd's from 0.51\,ms to 0.23\,ms and the mesh viewer's
from 0.35\,ms to 0.16\,ms (Release, idle screens). Lit dynamic meshes stay
inline so shadow bounds remain exact.\\
\textbf{Extends:} D-E14 (frame schema v3)
\end{decisionbox}

\subsection*{D-E17 --- Simulation lanes with retire-on-drain}
\begin{decisionbox}[title={D-E17 \quad 2026-09-22}]
\textbf{Decision:} IllumoGame generations run on up to eight isolated
simulation lanes: worker stores of \texttt{CSimWorkerGuest.wasm} reached
through the generic \texttt{LaneJob} and \texttt{JobLanes} services. Each lane
owns interleaved bands of eight chunk rows with a one-row halo and advances
them exactly with the unchanged sparse kernels;
\symbolref{SparseCellGrid}\texttt{::applyChunkPatches} keeps its frontier
journal sound across halo updates. The control store merges the owned changes
into the spare grid as one exact delta, launches the next generation from the
halos before merging, and publishes through the unchanged mirror path.
\symbolref{SimulationRunner}\texttt{::canBlock()} is false with lanes, so
edits, loads, ruleset changes, saves and exit retire the outstanding
generation instead of waiting for it. Lanes engage above 4\,ms serial
generations; elementary 1D rules (global source row), radii above 16 and lane
failures stay serial.\\
\textbf{Why:} Serial generations ran inside the frame's update call, so large
worlds froze frames; native production used a worker thread and 4--8 pool
workers.\\
\textbf{Consequences:} Frames no longer wait for generations: a
4,285-chunk dense world holds 60\,FPS at vsync instead of about 48, and
uncapped headless frame rates rise 3--4$\times$. Throughput is bounded by the control
store's merge (about 2\,ms per dense generation); measured uncapped TPS is
1.3--1.6$\times$ serial for dense worlds and remains below native production
(about 20--35\%). At most one generation is discarded when the world changes
while running. The whole-world \texttt{CSW1} worker remains for parity tests.\\
\textbf{Completes:} D-E13 (worker-store offload, milestone 6)
\end{decisionbox}

\subsection*{D-E18 --- Illumo::Content, the shared content layer}
\begin{decisionbox}[title={D-E18 \quad 2026-09-23}]
\textbf{Decision:} Add the optional static library \texttt{Illumo::Content}
(\texttt{Illumo/Include/Illumo/Content}, \texttt{Illumo/Source/Content}) above
the engine. It owns virtual paths, package manifests, \texttt{.ilpk} archives,
the virtual file tree, the \texttt{.ilsc} scene format and
\texttt{SceneInstance}. Core \texttt{Illumo} never includes
\texttt{<Illumo/Content/...>}; Content never depends on the WASM host or a
product; nlohmann stays private to its sources. Guests link the serial-safe
mirror \texttt{IllumoGuestContent} (paths, manifests, scene format, scene
instances); archives and the file tree stay host-side.\\
\textbf{Why:} IllEd, IllMeshViewer, IllumoGame and the runtime all needed the
same scene and package code. D-E10's condition for promoting the codec (a
second loader) was met.\\
\textbf{Consequences:} \texttt{SceneGraph} still never serializes itself;
Content builds graphs only through the public API. Engine tools reach the
file tree through the core interface \texttt{IFileTreeSource}
(\texttt{IllumoContext::fileTree}), which Content implements as
\texttt{VfsTreeSource}.\\
\textbf{Supersedes:} D-E10 (editor-owned codec)
\end{decisionbox}

\subsection*{D-E19 --- .ilsc format 2 is the one scene format}
\begin{decisionbox}[title={D-E19 \quad 2026-09-23}]
\textbf{Decision:} Every Illumo program reads and writes \texttt{.ilsc} format
2 (\texttt{format\_version} \texttt{[2, minor]}; newer minors are refused).
A document holds \texttt{metadata}, \texttt{settings} (world mode and one
environment: skybox, ambient colour, one directional sun), an \texttt{assets}
table (\texttt{mesh}, \texttt{texture}, \texttt{atlas}, cubemaps), nodes with
stable string ids, parents, transforms, tags and core components
(\texttt{primitive}, \texttt{mesh}, \texttt{sprite}, \texttt{light},
\texttt{camera}), namespaced components and \texttt{extensions} kept verbatim,
and an \texttt{editor} block. Core components are strict: unknown keys fail.
Asset references are package-relative (\texttt{meshes/tree.obj}) or explicit
absolute virtual paths (\texttt{/engine/...}, \texttt{/packages/<id>/...}).
Output is canonical: fixed key order and shortest round-trip floats.
\texttt{SceneInstance} is the live loader: it owns the document, graph,
attachments and id map, applies per-node edits incrementally, shares one unit
mesh per primitive kind, and falls back to placeholders for missing assets.\\
\textbf{Why:} Version 1 could only describe coloured primitives, could not
reference an asset, failed on any unknown kind, and only IllEd read it.
Package-relative references let a project directory, its packed
\texttt{.ilpk} and a mounted copy share one file unchanged.\\
\textbf{Consequences:} Clean break: version 1 files are not read. Loads
collect references (\texttt{collectSceneFetches}), fetch them, then
instantiate, so \texttt{IAssetSource} stays synchronous. IllumoGame's
\texttt{render3dTest} and IllMeshViewer scenes use the same path.
\end{decisionbox}

\subsection*{D-E20 --- .ilpk packages are a bounded ZIP subset}
\begin{decisionbox}[title={D-E20 \quad 2026-09-23}]
\textbf{Decision:} A package is a directory or an \texttt{.ilpk} archive: a
ZIP subset with stored and deflate entries only (no zip64, multi-disk,
encryption or symlinks). The reader rejects unsafe or duplicate names,
file/directory collisions and local headers that disagree with the central
directory; it bounds entry count (65,535), directory bytes (16\,MiB),
per-entry (256\,MiB) and total (1\,GiB) inflated size and the compression
ratio (1,024), and checks CRC-32 on every full read. The host writer
(\texttt{PackageArchiveWriter}, the \texttt{IllumoPack} tool and the Pack file
action) deflates through the vendored \texttt{stb\_image\_write} compressor and
stores already-compressed images as they are.\\
\textbf{Why:} Mods and expansions need one shippable file that standard tools
can inspect, without a new dependency and without trusting the archive.\\
\textbf{Consequences:} Inflate uses \texttt{stb\_image}'s decoder with an
exact output bound. CMake's own zip output is a test fixture, so the reader is
checked against an independent writer.
\end{decisionbox}

\subsection*{D-E21 --- Packages mount on one host virtual file tree}
\begin{decisionbox}[title={D-E21 \quad 2026-09-23}]
\textbf{Decision:} \texttt{IllumoRuntime} owns a \texttt{VirtualFileSystem}.
It mounts \texttt{/engine} (runtime assets), \texttt{/app} (the launched
package plus overlays), \texttt{/packages/<id>} (every package found in
\texttt{packages/} beside the runtime and each \texttt{--mount}), and a
writable \texttt{/project} only with \texttt{--project <dir>}. Mount tables are
immutable snapshots swapped under a mutex; open files keep their backend
alive after an unmount. Overlays target only \texttt{/app}: priority, then
dependency order, then id; directories union and the top file wins; there are
no whiteouts. Directory lookups are case-exact on every file system, and a
path that escapes its root through a link does not exist. \texttt{stat} names
the supplying package but never a host path.\\
\textbf{Why:} Content must be browsable at runtime and in the console, and
several packages must coexist for mods and expansions without any package
choosing its own mount point.\\
\textbf{Consequences:} The host \texttt{vfs} console command lists, stats and
reads the tree, and debug builds add the \texttt{files} browser. Writes go only
to \texttt{/project}; \texttt{*.wasm} and \texttt{illumo.json} are refused
there. A mod cannot delete a base file.
\end{decisionbox}

\subsection*{D-E22 --- illumo.json is the one package manifest}
\begin{decisionbox}[title={D-E22 \quad 2026-09-23}]
\textbf{Decision:} Every package carries \texttt{illumo.json}
(\texttt{format ilpk}, \texttt{format\_version} 1, \texttt{id},
\texttt{version}, \texttt{title}, \texttt{kind} \texttt{app}, \texttt{content}
or \texttt{mod}, \texttt{targets}, \texttt{dependencies}, \texttt{overlays},
and an \texttt{app} or \texttt{mod} section). The \texttt{app} section absorbs
the former \texttt{app.json} fields and their clamping to host ceilings. A
\texttt{kind} whose sections do not match is refused; dependency order is
Kahn's algorithm and cycles are refused.\\
\textbf{Why:} Apps, content and mods are all packages; two manifest formats
would have diverged.\\
\textbf{Consequences:} \texttt{app.json} and \texttt{AppManifest} are removed;
staging emits \texttt{illumo.json}; \texttt{--app} and \texttt{--package}
accept a directory or an \texttt{.ilpk}.\\
\textbf{Refines:} D-E14
\end{decisionbox}

\subsection*{D-E23 --- File protocol v2 and the ProjectFiles capability}
\begin{decisionbox}[title={D-E23 \quad 2026-09-23}]
\textbf{Decision:} Guest file requests gain a \texttt{Mounted} area addressed
by absolute virtual paths, and the Package area becomes an alias for
\texttt{/app}. New actions: \texttt{List} (paged), \texttt{Stat},
\texttt{Import} (copy a dialog grant into \texttt{/project}) and \texttt{Pack}
(\texttt{/project} to an \texttt{.ilpk} on a dedicated host worker). Mounted
reads use blocks of up to 1\,MiB, and a guest may keep 16 file tasks in
flight. Writing a mounted path, Import and Pack need the new capability
\texttt{ProjectFiles} (bit 9), which the host offers only while a project is
mounted; the project quota is tracked incrementally.\\
\textbf{Why:} Guests could read only their own flat package directory; the
editor needs to browse every mount and write a project without learning host
paths.\\
\textbf{Consequences:} Reading mounted paths needs only \texttt{Assets}.
A pack never replaces a currently mounted archive: the replace fails and the
guest sees \texttt{IoError} (there is no up-front check). The guest SDK adds
\texttt{GuestFileTree}. Every new request shape has decoder and deny tests.
\end{decisionbox}

\subsection*{D-E24 --- Guest asset cache over the file tree}
\begin{decisionbox}[title={D-E24 \quad 2026-09-23}]
\textbf{Decision:} \texttt{GuestVfsAssets} replaces \texttt{GuestPackageAssets}
as the guest \texttt{IAssetSource}. \texttt{packageAssets()} remains a pinned
bootstrap preload; loads fetch sets of virtual paths on demand and hold them
until released; unpinned, unheld bytes are evicted least recently used past a
256\,MiB budget; \texttt{/local/<name>} holds bytes the app supplies itself.
\texttt{GuestSceneFetches} (\texttt{IllumoGuestContent}) adds the second stage
that fetches the material libraries an OBJ declares.\\
\textbf{Why:} Asset bytes must be readable synchronously when a scene
instantiates, but a guest cannot preload every package it might open.\\
\textbf{Consequences:} IllEd and IllMeshViewer share one fetch pipeline; a
missing material library is tolerated, a missing asset is reported and drawn
as a placeholder.
\end{decisionbox}

\subsection*{D-E25 --- IllEd editing model}
\begin{decisionbox}[title={D-E25 \quad 2026-09-23}]
\textbf{Decision:} IllEd edits a format 2 document over \texttt{SceneInstance}.
Each command is a patch of node records (before, after, parent, sibling
position); drags merge into one command; history is capped at 512 commands
and 64\,MiB and cleared on load; dirty means the history cursor differs from
the saved one. Selection is a set with a primary node. One shortcut table
drives menus and keys. New nodes go under the root unless ``create as child''
is chosen. Picking uses \texttt{raycastCandidates} in both world modes.
Clipboard fragments (\texttt{illed.fragment}, at most 4\,MiB) remap ids on
paste. With \texttt{--project}, the asset browser places files from the tree,
and Import, Save to Project and Pack Project write through D-E23.\\
\textbf{Why:} The prototype had no undo, one selected node, world-axis
translation only, a read-only inspector and no assets; camera motion marked
documents dirty.\\
\textbf{Consequences:} The module is split into document, history, selection,
shortcut, gizmo, inspector, clipboard, asset and project units. Console
commands \texttt{scene\_select}, \texttt{scene\_place} and
\texttt{scene\_save\_project} script the same flows for tests.
\end{decisionbox}

\subsection*{D-E26 --- SceneGraph sibling order}
\begin{decisionbox}[title={D-E26 \quad 2026-09-23}]
\textbf{Decision:} Add one narrow ordering call to SceneGraph v2:
\texttt{setParent(node, parent, insertBefore)} places a node immediately
before a sibling, or last when \texttt{insertBefore} is null. An unchanged
position is a no-op, and invalid or cyclic requests fail without mutating.\\
\textbf{Why:} Scene documents and the editor hierarchy are ordered, and undo
must restore a node to its exact position.\\
\textbf{Consequences:} \texttt{SceneInstance} and IllEd reorder, paste and undo
in place. No other part of the v2 contract changes.\\
\textbf{Refines:} D-E12
\end{decisionbox}

\subsection*{D-UI6 --- GuiTextEdit and GuiFileTree}
\begin{decisionbox}[title={D-UI6 \quad 2026-09-23}]
\textbf{Decision:} Two primitive-composed helpers join \texttt{Illumo/Gui},
native and guest. \texttt{GuiTextEdit} is a UTF-8 line editor (caret,
selection, clipboard) fed by \texttt{InputManager::getCharQueue} and drawn by
\texttt{GuiKit::drawTextField}. \texttt{GuiFileTree} keeps expansion state
and the listings received so far, and flattens them without recursion into
rows drawn with \texttt{GuiKit::drawTreeRow}. There is no retained widget
tree.\\
\textbf{Why:} IllEd's inspector needed typed fields, and the asset browser and
the debug file browser needed one tree view over asynchronous listings.\\
\textbf{Consequences:} The IllEd inspector and asset browser use them, and so
does DebugModule's \texttt{files} command, which browses
\texttt{IllumoContext::fileTree} with the keyboard and wheel only, so every
input it takes is consumed before the product sees it.\\
\textbf{Extends:} D-UI1 (menu shell)
\end{decisionbox}

\subsection*{D-UI7 --- Plain tool look and detachable panels}
\begin{decisionbox}[title={D-UI7 \quad 2026-09-23}]
\textbf{Decision:} IllEd and IllMeshViewer share a plain tool look
(\texttt{GuiToolStyle}: a neutral palette, flat panels, title bars, splitters,
rows, buttons, toggles, sliders, text fields, menus and status bars) and one
\texttt{GuiPanelDock} per app window. The dock stacks panels in left and right
columns with draggable splitters; each title bar can hide a panel or pop it
out into its own window, and dragging a title past the window edge tears it
off, as the console does. A detached panel docks back from its title bar, its
window's close box or View \textgreater{} Reset Layout. Panels are content
renderers: they draw into a \texttt{GuiPanelPlacement} (content rectangle
plus surface) and read their pointer through \texttt{GuiPanelPointer}, which
samples the main window or the panel's window and honours presses other
chrome took. The layout is saved in the \texttt{panelLayout} setting and
detached panels reopen on the next start.\\
\textbf{Why:} The owner asked for the editor and viewer interfaces to be
redone with panels that break off into separate windows, like the developer
console, in a plain tool look.\\
\textbf{Consequences:} IllEd's Hierarchy, Assets, Tools and Inspector panels
and the viewer's Info and Display panels detach; an asset dragged from a
detached Assets window places where it is dropped in the viewport. The chrome
keeps fixed tool metrics while \texttt{fontSize} scales panel content. There
is still no retained widget tree.\\
\textbf{Extends:} D-UI4 (plain console palette), D-UI5 (console pop-out)
\end{decisionbox}

\subsection*{D-E27 --- Panel surfaces}
\begin{decisionbox}[title={D-E27 \quad 2026-09-23}]
\textbf{Decision:} \texttt{IPanelSurfaces}, published as
\texttt{IllumoContext::panelSurfaces}, gives products extra top-level windows.
Surface 0 is the main window; others are opened, titled and closed by id, and
report size, screen origin, pointer and window events. Keys stay in the one
\texttt{InputManager} queue and \texttt{focused()} names their window. Guests
get it through the \texttt{Windows} capability (bit 10), which the host grants
only when it can present windows, the Window service (17) and input v2;
\texttt{WasmPanelWindows} owns the host windows and closes them on failure and
exit. Products must work fully docked without it.\\
\textbf{Why:} Product modules are shared by the guest and the native test
oracles, so they need a seam free of guest types; \texttt{FakePanelSurfaces}
drives the same code in tests.\\
\textbf{Consequences:} Linux, \texttt{--capture}, \texttt{--bench-*} and
headless runs keep every panel docked with pop-out disabled.\\
\textbf{Extends:} D-E23 (capabilities)
\end{decisionbox}

\subsection*{D-E28 --- Sound effects}
\begin{decisionbox}[title={D-E28 \quad 2026-09-24}]
\textbf{Decision:} \texttt{IAudio}, published as \texttt{IllumoContext::audio},
registers decoded clips as sounds (at most 256) and plays them by
generational handle, one voice per play (at most 32, the oldest replaced),
with volume, pan, pitch, \texttt{stopAll} and a master volume. The native
implementation, \texttt{AudioDevice}, and the \texttt{AudioDecoder} for WAV,
FLAC and MP3 wrap vendored miniaudio 0.11.25, which stays private to
\texttt{Illumo/Source/Audio}. Guests get \texttt{IAudio} through the
\texttt{Audio} capability (bit 11), granted only when the host has an output
and never for \texttt{--capture} or \texttt{--bench-*}, and the Audio service
(18): \texttt{GuestAudioRequest} version 1 with Create, Destroy, Play, StopAll
and SetVolume. Requests complete in the same exchange without a payload and
the guest queue discards their completions.\\
\textbf{Why:} Guests decode their own package files, as they decode textures,
so the host never parses untrusted encoded audio; only range-checked float
samples cross the ABI. The decoder uses miniaudio's memory entry points so
no file-system import reaches a guest. The public seam keeps products free of
library and guest types, as D-E27 did for windows.\\
\textbf{Consequences:} A clip holds at most 3\,Mi samples and a guest at most
64\,MiB of registered samples; beyond that, and for unknown ids, requests are
rejected without retiring the guest. The host releases a guest's sounds and
restores the master volume when it retires. Products must also work
silently. CSim plays seven cues scaled by its \texttt{soundVolume} setting.\\
\textbf{Extends:} D-E27 (capability pattern)
\end{decisionbox}

\subsection*{D-E29 --- Lanes for every rule from 1\,ms}
\begin{decisionbox}[title={D-E29 \quad 2026-09-25}]
\textbf{Decision:} The guest runner moves generations to simulation lanes once
serial generations exceed 1\,ms (previously 4\,ms) and returns them to the game
store when all lanes together work under 0.5\,ms (previously 1\,ms).
Elementary 1D rules also take lanes: their partition runs along chunk columns,
the coordinator finds the global source row (the highest row with a counted
cell and its counted extent) and sends it with each Advance, each lane writes
only its own columns of the next row, and every lane reports the source row of
its owned columns so the coordinator combines the next one without the merged
grid. The lane protocol becomes CSL1 version 2: SyncBegin carries the axis;
Advance and Advanced carry the row; decoders reject version 1, an inverted row
extent, and an axis or row that does not match the rule.\\
\textbf{Why:} \texttt{IllumoGame.Wasm.PackageBench} shows that at about 1\,ms
a lane generation costs the game store only its merge (0.30--0.41\,ms on 1\,ms
worlds, median frame 1.36\,ms to 0.55\,ms), so the old threshold left
frames carrying work lanes could take. Elementary 1D was excluded only because
a row band cannot see the global source row; supplying the row removes the
reason, and each lane then also copies only its own columns of the growing
history.\\
\textbf{Consequences:} Uncapped peak TPS on light worlds drops by roughly a
fifth (the round trip, about 1.3\,ms), which normal speeds never reach; frame
time falls instead. \texttt{IllumoGame.Wasm.LaneParity} now requires lanes for
every catalog rule and adds wide elementary rows on infinite and wrapping
worlds with an edit; \texttt{IllumoGame.Wasm.LaneProtocol} covers version 2.
Two light benchmark worlds (\texttt{bench-dense16}, \texttt{bench-sparse64})
cover the threshold region.\\
\textbf{Supersedes:} the 4\,ms entry and the elementary exclusion of D-E17
\end{decisionbox}

\subsection*{D-R27 --- Surface frames and offscreen readback}
\begin{decisionbox}[title={D-R27 \quad 2026-09-23}]
\textbf{Decision:} Frame schema v5 adds a trailing section of at most eight
surfaces, each a UI-only batch list with a revision, or \texttt{same} when
nothing it draws changed. The host replays a changed surface through
\texttt{Renderer::renderOffscreen} into a per-window target and queues an
asynchronous readback (\texttt{IBackend::requestFramebufferReadback} and
\texttt{takeFramebufferReadback}, two fenced pixel-pack buffers per stream);
the newest completed copy is presented by the window's \texttt{PixelWindow},
one frame late.\\
\textbf{Why:} Detached panels must draw exactly as docked ones (sprites,
ellipses, text) without a second GL context or a CPU rasterizer.\\
\textbf{Consequences:} An unchanged panel costs neither a replay nor a
readback. Versions 1--4 remain valid, and the frame quotas span the main frame
and every surface.\\
\textbf{Extends:} D-UI5 (no second GL context), D-R25
\end{decisionbox}

\subsection*{D-UI8 --- Liquid, bouncy menu motion}
\begin{decisionbox}[title={D-UI8 \quad 2026-09-23}]
\textbf{Decision:} The glass menus move like liquid. A travelling selection
is a drop whose head and tail are two springs
(\texttt{GuiMotion::kLiquidHead}, \texttt{kLiquidTail}) owned by
\symbolref{GuiMenuAnimator}: from rest it leaves the old row, mid-flight it
keeps its momentum, and \texttt{GuiSelectionSpan::squash} reports the jelly
response from the closing speed between head and tail.
\texttt{GuiKit::drawLiquidSelection} draws it as non-overlapping slices across
the travel: the head cell keeps full width, the drop necks behind it and
tapers toward the tail like a teardrop, and at rest it is exactly a rounded
rectangle. Panels and rows drop in on springs
(\texttt{GuiEasing::springStep}), presses and stepped values wobble
(\texttt{GuiEasing::wobble}, \texttt{pressWobble}, \texttt{valueWobble}), and
screens tune every spring from shared \texttt{GuiMotion} presets rather than
literal frequencies. \texttt{GuiKit::drawSplash} adds a ripple and teardrop
droplets for a press. Every glass panel swivels toward the pointer through
\texttt{GuiPanelTilt} (a \texttt{GuiMotion::kSway} spring pair): its layers
shift by depth and \texttt{GuiGlassStyle::tiltX}/\texttt{tiltY} slide the
shadow away, draw a glare that follows the pointer and light the facing
edges; each screen shifts its layout origin by the body layer so hit testing
follows the drawing. The simulator's corner EDIT/NORMAL label becomes a
product-drawn \texttt{ModeBadge} in the same glass (superseding D-OWN1's
\symbolref{SplashText} member; the engine class stays).\\
\textbf{Why:} The owner asked for menu animation that feels bouncier, more
alive, and more like liquid and teardrops.\\
\textbf{Consequences:} \texttt{beginSelectionTravel} takes the row left and
the row targeted; a selection changed without a travel reads as its own row.
Hit areas still never follow animated offsets beyond the entrance, reduced
motion snaps every new spring and wobble, and the rounded
\symbolref{GuiDialog} shares the same drop horizontally while its flat style
is unchanged.\\
\textbf{Extends:} D-UI1 (menu shell), D-R26 (living-glass chrome)
\end{decisionbox}

\subsection*{D-UI9 --- Variable-weight type for CSim}
\begin{decisionbox}[title={D-UI9 \quad 2026-09-23}]
\textbf{Decision:} CSim sets its type in Kikuta (SIL OFL 1.1), a
variable-weight face shipped under \texttt{Assets/Fonts/Kikuta}. The host font
service keeps its closed catalog and adds the name
\texttt{kikuta:<weight>[:<glyphs>]}: FreeType instances the \texttt{wght} axis
(\texttt{FontFaceOptions}), rasterizes the printable ASCII Kikuta lacks from
Space Mono, and can restrict an atlas to up to 32 glyphs. Weight animates
without a live rasterizer: \texttt{FontWeightRamp} holds a few rasterized
weights and maps any weight to the two around it, and
\texttt{TextPrimitive::heavyFont} overlays the heavier at the fraction between
them with interpolated advances. \texttt{CSimTypeface} installs Kikuta 400 as
the default font and a 400--1000 UI ramp during package bootstrap;
\texttt{GuiKit::drawEmphasizedText} and its centered and measuring companions
turn a row or button's emphasis spring into weight, and the CSIM title animates
its own four-glyph ramp (thin while falling, heavy on impact, a rolling weight
swell and a pull toward the pointer). \texttt{TextPrimitive::stretchX} and
\texttt{stretchY} squash and stretch a run about its left edge and baseline;
the title letters use them with weight for random one-at-a-time poses (flex,
slim, hop) and a hop wave when the word is clicked.\\
\textbf{Why:} The owner supplied the face and asked for hover emphasis to
thicken labels and for play with the title's weight.\\
\textbf{Consequences:} The guest ABI and \texttt{GuestFontRequest} are
unchanged. Products without an installed ramp draw emphasized text plainly, so
IllEd and the viewer are unaffected. Each ramp sample is one font request
against the guest's 32-font budget (CSim uses at most 20). Kikuta's B and S
resemble 8 and 5, which rule strings such as B3/S23 inherit.\\
\textbf{Extends:} D-UI8 (menu motion), D-R26 (living-glass chrome)
\end{decisionbox}

\subsection*{D-UI10 --- CSim's software pointer}
\begin{decisionbox}[title={D-UI10 \quad 2026-09-24}]
\textbf{Decision:} CSim draws its own mouse pointer (\texttt{SoftwareCursor})
and hides the system cursor while it does. Display wire version 2 appends
\texttt{hideSystemCursor} to \texttt{GuestDisplayState}; the host applies it
with GLFW's hidden-cursor mode on the main window only
(\texttt{IRenderWindow::setSystemCursorHidden}), still decodes version 1 and
answers each request in its own version, and a version 1 request leaves the
cursor alone. \texttt{WasmGameServices::cancel} restores the system cursor, so
a stopped or failed guest never leaves it hidden.
\texttt{GuestModuleApplication} gains \texttt{updateOverlay} and
\texttt{dispatchOverlay} for product overlays that outlive module transitions
and draw above them. The pointer's tip sits exactly on the pointer; its body
leans against sideways motion on a spring, a faint holographic afterimage
follows movement, and a press squishes it and throws a splash. It takes the menu
glass palette: a dark halo, a rim lit cyan at the tip and violet at the tail,
a deep glass body, and a soft glow around the edge.\\
\textbf{Why:} The owner asked for a software cursor with sway and
movement-driven animation.\\
\textbf{Consequences:} The system cursor returns while the host console is
open (it draws above the game) and outside the window's content area.
The Software cursor toggle in F1 settings (the persisted
\texttt{softwareCursor} option, also \texttt{toggle softwareCursor}) turns it
off.
IllEd and the viewer never set the field. The drawn pointer follows the last
input snapshot, so it can trail the hardware pointer by a frame.\\
\textbf{Extends:} D-UI8 (menu motion), D-E23 (capabilities)
\end{decisionbox}
\subsection*{D-PL1 --- Retire the macOS scaffold (2026-09-19)}
macOS is not a current product target. Remove the first-party Apple platform
files, Cocoa framework/source selection, and dedicated scaffold package map.
CMake rejects Apple targets explicitly. Windows remains supported; Linux keeps
its existing source-repaired, native-validation-pending status. Vendored
dependency portability code and historical design records retain their provenance.
A future macOS port requires a new scope and native validation.
